\documentclass[11pt]{article}

\usepackage{fontspec}
\defaultfontfeatures{Renderer=Harfbuzz,Ligatures=TeX}
\setmainfont{Noto Sans}
\setsansfont{Noto Sans}
\setmonofont{Noto Sans Mono}
\usepackage[margin=1in]{geometry}
\usepackage{amsmath,amssymb}
\usepackage{booktabs}
\usepackage{array}
\usepackage{longtable}
\usepackage{tabularx}
\usepackage{graphicx}
\usepackage{caption}
\usepackage{enumitem}
\usepackage{ragged2e}
\usepackage{hyperref}
\usepackage{url}
\usepackage{polyglossia}
\setdefaultlanguage{ukrainian}
\newfontfamily\cyrillicfont[Script=Cyrillic]{Noto Sans}
\newcolumntype{P}[1]{>{\raggedright\arraybackslash}p{#1}}
\captionsetup{font=small,labelfont=bf}
\setlength{\emergencystretch}{6em}
\hypersetup{
  unicode=true,
  colorlinks=true,
  linkcolor=blue,
  citecolor=blue,
  urlcolor=blue,
  pdflang={uk},
  pdftitle={Cosmo-Local Credit (CLC) White Paper v0.8},
  pdfauthor={William O. Ruddick and Mohamed Sohail}
}

\title{Cosmo-Local Credit (CLC)\\Мережа для маршрутизації кредиту, координації зобов'язань і фінансування здорової космо-локальної економіки}
\author{William O. Ruddick \\ Mohamed Sohail \\ Grassroots Economics Foundation \\ \texttt{info@grassecon.org}}
\date{White Paper v0.8 translation: 10 October 2026}

\begin{document}
\maketitle
\tableofcontents
\newpage
\sloppy
\RaggedRight

\textbf{Про цей переклад.} Це переклад White Paper v0.8. Англійський вихідний текст опубліковано 30 вересня 2026 року. Цей переклад опубліковано 10 жовтня 2026 року. Англійська версія є вихідним текстом. \href{https://docs.cosmolocal.credit/white-paper}{Прочитати англійський вихідний текст}.

\textbf{Версія:} 0.8

\textbf{Дата:} 30 вересня 2026 року

\textbf{Автори:} William O. Ruddick і Mohamed Sohail

\textbf{Контакт:} \href{mailto:info@grassecon.org}{info@grassecon.org}

\textbf{Завантажити:} \href{https://docs.cosmolocal.credit/white-paper/Cosmo-Local-Credit-CLC-White-Paper-v8-uk.pdf}{CLC White Paper v0.8 (PDF українською)}

\textbf{Публікація:} остаточне публічне видання з номером версії. Цей документ не є інвестиційною порадою й не становить пропозиції продати або заклику придбати цінний папір чи фінансовий інструмент. Подальші версії можуть змінювати описані тут моделі.

\section*{Як читати цей документ}
\addcontentsline{toc}{section}{Як читати цей документ}

Ця Біла книга поєднує чотири види матеріалів. Статус застосовується до всього документа, навіть якщо окремий розділ не повторює його.

\begin{longtable}{P{0.30\linewidth} P{0.62\linewidth}}
\toprule
Статус & Значення \\
\midrule
\endhead
\textbf{Поточне} & Поведінка, перевірена в CLC App або публічних контрактах Protocol v1.1.0. \href{https://docs.cosmolocal.credit/uk/protocol/overview}{Довідка про Протокол} є фактичним джерелом щодо цих контрактів. \\
\textbf{Залежить від розгортання} & Можливість, що існує лише там, де її вмикає конкретний Пул, постачальник, юрисдикція, орган управління або опублікована політика. \\
\textbf{Історичне} & Дані або функції попереднього сервісу Sarafu Network. Це не поточне використання CLC і не доказ того, що користувача, актив, запис або зобов'язання було перенесено. \\
\textbf{Пропоноване} & Дизайн, економічна модель, механізм управління або пункт дорожньої карти, який не представлено як розгорнутий. \\
\bottomrule
\end{longtable}

Поточний Protocol v1.1.0 охоплює пряме виконання через \texttt{SwapPool}, необов'язкове кураторство й оцінювання, верхні межі балансу токенів Пулу, комісії Пулу, додаткову протокольну комісію та \texttt{SwapRouter}, який лише надає котирування. Багатокрокове виконання, маршрутизація через HTLC або ескроу, пакетний взаємозалік, спільне страхування, пропонований токен управління CLC, пропонований мережевий Пул CLC, механізми кредитів на комісії та програми мережевої ліквідності є пропонованими або залежать від розгортання.

\textbf{Цільова аудиторія:} розробники й Куратори Пулів, інженери та аудитори протоколу, фінансувальники ліквідності й мандатів, розробники механізмів управління та фахівці з права або відповідності вимогам.

\section*{Анотація}
\addcontentsline{toc}{section}{Анотація}

Cosmo-Local Credit --- це сімейство продуктів, побудоване навколо \textbf{Протоколу об'єднання зобов'язань (CPP)}. CPP описує, як незалежно керовані Пули зобов'язань можуть відбирати зобов'язання, що підлягають погашенню, публікувати методи обмінного курсу й ліміти, зберігати запаси та забезпечувати відповідальний обмін.

\textbf{Пул зобов'язань} --- це керована домовленість. Поточний \textbf{\texttt{SwapPool}} у Protocol v1.1.0 є однією контрактною складовою: сховищем токенів і механізмом прямого обміну. Розгортання може під'єднати цей контракт до реєстрів, модулів котирування, верхніх меж балансу токенів Пулу, політик комісій та інших сервісів.

Космо-локальний підхід означає глобально поширювати відкриті стандарти, програмне забезпечення й знання, залишаючи випуск, виконання, управління та реальну відповідальність на місцевому рівні. Емітент залишається відповідальним за свій ваучер. Куратор Пулу --- за правила Пулу й кожну прямо взяту гарантію. Ані CLC App, ані GEF автоматично не гарантують ваучер, Пул, вартість, маршрут, позицію ліквідності чи результат.

\section*{Поточне розгортання}
\addcontentsline{toc}{section}{Поточне розгортання}

GEF керує CLC App на \href{https://cosmolocal.credit}{cosmolocal.credit}. Він надає доступ до облікового запису й гаманця, публічний каталог Ринку та підтримувані інтерфейси для токенів, ваучерів, Пропозицій, Пулів зобов'язань, переказів і прямих обмінів через Пул. Доступність залежить від інтерфейсу, контрактів, запасів, налаштованих лімітів і комісій, стану мережі, відповідності вимогам, постачальників, юрисдикції та опублікованих умов емітента або Пулу.

Керування інтерфейсом не робить GEF автоматично емітентом, Куратором Пулу, зберігачем, кредитором, позичальником, брокером, гарантом, стороною, що погашає ваучер, радником, страховиком або стороною зобов'язань між учасниками. Використання застосунку регулюють \href{https://docs.cosmolocal.credit/uk/governance/terms}{Умови користування}.

Protocol v1.1.0 відокремлює відповідальні ролі від технічного контролю. Куратор Пулу, власник \texttt{SwapPool}, адміністратор проксі, контролер залежності, модератор каталогу й отримувач комісії можуть бути різними сторонами. Спільну термінологію наведено на сторінці \href{https://docs.cosmolocal.credit/uk/introduction/concepts}{Поняття й термінологія}.

\section*{Історичне походження}
\addcontentsline{toc}{section}{Історичне походження}

Історична Sarafu Network передувала CLC App і надала практичні дані про громадські валюти та об'єднання зобов'язань. Перехід сервісу до Cosmo-Local Credit не переніс автоматично кожен історичний обліковий запис, гаманець, токен, ваучер, Пул, баланс, звіт, учасника або зобов'язання.

Історичні показники в цій Білій книзі походять із датованого, внутрішньо узгодженого набору даних Sarafu. Це не поточні показники використання CLC. Джерело й період вимірювання наведено в \href{https://docs.cosmolocal.credit/uk/white-paper/executive-summary}{Короткому викладі}.

\section*{Модель дизайну}
\addcontentsline{toc}{section}{Модель дизайну}

CPP виконує чотири широкі функції:

\begin{itemize}[leftmargin=*]
\item \textbf{Кураторство:} визначення ваучерів або інших активів, які допускаються.
\item \textbf{Оцінювання:} публікація методу, що створює обмінні курси або котирування Пулу.
\item \textbf{Обмеження:} застосування поточних верхніх меж балансу токенів Пулу або інших окремо реалізованих засобів контролю ризику.
\item \textbf{Обмін:} зберігання запасів і виконання обмінів через Пул за розкритими комісіями й межами.
\end{itemize}
Розгортання може додати маршрутизацію, моніторинг, підтримку ліквідності, гарантії, резерви, страхування або сервіси управління лише за окремо опублікованими політиками. Розміщення, кураторство, ліміт, резерв або блокчейн-запис самі по собі не доводять спроможність емітента, виконання, схвалення, гарантію чи страхування.

\section*{Цінності й принципи дизайну}
\addcontentsline{toc}{section}{Цінності й принципи дизайну}

\begin{itemize}[leftmargin=*]
\item \textbf{Турбота про людей:} підтримувати індивідуальне й колективне добробут.
\item \textbf{Турбота про довкілля:} зменшувати екологічну шкоду й підтримувати відновлювальні практики.
\item \textbf{Справедливість:} сприяти безпечному й рівноправному доступу до ресурсів, знань і турботи.
\item \textbf{Взаємність:} ділити ризик, витрати, відповідальність і надлишок за розкритими правилами.
\item \textbf{Недопущення домінування:} протидіяти зосередженню контролю над людьми, даними, фінансами, знаннями й матеріальними ресурсами.
\item \textbf{Стійкість:} допомагати громадам готуватися до економічних, політичних і кліматичних потрясінь та адаптуватися до них.
\item \textbf{Місцева відповідальність:} залишати реальне виконання й юридичну відповідальність за визначеними сторонами.
\item \textbf{Можливість відокремлення:} давати сумісним громадам і розгортанням змогу залишати спільні реєстри або сервіси без стирання чинних балансів чи зобов'язань.
\end{itemize}
Цифровий рівень --- це спільна пам'ять та інфраструктура координації. Він не замінює відносини, виконання емітентом, місцеве управління чи юридичну відповідальність.

\section*{Поточні й пропоновані процеси}
\addcontentsline{toc}{section}{Поточні й пропоновані процеси}

\subsection*{Поточний прямий обмін через Пул}
\addcontentsline{toc}{subsection}{Поточний прямий обмін через Пул}

\begin{enumerate}[leftmargin=*]
\item Власник перевіряє ваучер, його емітента, умови й пов'язані Пропозиції.
\item Пул допускає підтримувані активи й публікує свою конфігурацію та відповідальних осіб.
\item Власник отримує котирування й авторизує прямий обмін через Пул.
\item \texttt{SwapPool} здійснює розрахунок за обміном у блокчейні та обліковує комісії Пулу й протоколу.
\item Якщо пізніше власник вибирає в застосунку \textbf{«Погасити»}, застосунок готує переказ власнику токена як пред'явлення для погашення.
\item Емітент окремо виконує обіцяну дію й записує списання, щоб виконані одиниці не можна було використати повторно.
\end{enumerate}
Отримання запасу з Пулу є обміном. Воно є погашенням лише тоді, коли Пул окремо взяв на себе роль емітента й виконує відповідні умови ваучера.

\subsection*{Ширший пропонований дизайн}
\addcontentsline{toc}{subsection}{Ширший пропонований дизайн}

Пропонований дизайн досліджує маршрутизацію виконання, взаємозалік, спільні сервіси, програми мережевої ліквідності, страхові політики й активи управління. Кожен із цих механізмів потребує реалізації, відповідальних сторін, ухвалених правил, фінансування, контролю й публічних розкриттів, перш ніж його можна буде застосувати.

Пропонована модель комісій може використовувати частку мережі з комісій Пулу та окремі комісії за маршрутизацію або сервіси. Ця пропозиція відрізняється від Protocol v1.1.0, де необов'язкова протокольна комісія додається до комісії Пулу.

Поточний \texttt{SwapPool} не надає прихильникам або постачальникам ліквідності автоматичних часток Пулу чи прав на повернення, виведення, винагороду або управління. Кожна поточна чи майбутня програма має окремо розкривати вид передавання, відповідальну сторону, контроль, права на відновлення або виведення, збитки, винагороди, комісії та права управління.

\section*{Навігація для читача}
\addcontentsline{toc}{section}{Навігація для читача}

\begin{itemize}[leftmargin=*]
\item Поточна термінологія й життєвий цикл дій: \href{https://docs.cosmolocal.credit/uk/introduction/concepts}{Поняття й термінологія}
\item Перевірені розгорнуті контракти: \href{https://docs.cosmolocal.credit/uk/protocol/overview}{Protocol v1.1.0}
\item Історичні дані: \href{https://docs.cosmolocal.credit/uk/white-paper/executive-summary}{Короткий виклад}
\item Федерація та взаємодія: розділи 5, 8 і 11
\item Пропоновані межі гарантій і страхування: розділи 10 і 11
\item Пропоновані моделі ліквідності й комісій: розділи 7, 9 і 12
\item Пропонована система вимірювання: розділ 14 і додатки A--C
\item Пропонована дорожня карта: розділ 15
\item Правові питання й відповідність вимогам: розділ 17
\end{itemize}
\href{https://docs.cosmolocal.credit/uk/white-paper/archive}{Попередні версії Білої книги} зберігаються лише як чітко позначені застарілі історичні публікації.

\section*{Короткий виклад}
\addcontentsline{toc}{section}{Короткий виклад}

Останні \textbf{Протокол об'єднання зобов'язань (CPP)} описує, як самостійно керувана Пули зобов'язань може регулювати погашувальні зобов'язання, публікувати методи та обмеження обмінного курсу, вести інвентаризацію та забезпечити відповідальний обмін. Емітенти залишаються відповідальними за свої зобов'язання з ваучера. Куратори Пулів залишаються відповідальними за правила Пул та будь-які гарантії, які вони прямо беруть на себе. Ні CLC App, ні GEF не є загальним гарантом.

Protocol v1.1.0 забезпечує поточні елементи прямих обмінів: \texttt{GiftableToken}, \texttt{SwapPool}, факультативні реєстри та модулі оцінки, обмеження балансу токенів Пул, плати Пул, додаткова плата протоколу \texttt{SwapRouter}. Подійні контракти не зафіксують реалізацію в реальному світі, не створюють автоматичні акції постачальника ліквідності, не виконують багатопропускних маршрутів або не забезпечують універсальні резерви, гарантії або страхування.

Історичне \textbf{Sarafu Network} використовується концепція об'єднання зобов'язань у валютах спільноти, заощаджувальних групах, системах взаємопомоги та виробничих зобов'язаннях. Його послуга пізніше перейшла на CLC App. Це не було відображенням того, що кожен історичний рахунок, гаманець, токен, ваучер, фонд, баланс, звіт, учасник або зобов'язання мігрували.

\subsection*{Історичний хвіст Sarafu Network --- діяльність Celo з 5 липня 2023 року до 20 липня 2025 року}
\addcontentsline{toc}{subsection}{Історичний хвіст Sarafu Network --- діяльність Celo з 5 липня 2023 року до 20 липня 2025 року}

\textbf{Джерело:} \href{https://dune.com/grassrootseconomics/sarafu-network}{Дашборд Dune Analytics}

\begin{itemize}[leftmargin=*]
\item 26 367 користувачів
\item 285 197 обмінів середніх осіб
\item 188 унікальний активний Пули зобов'язань
\item 745 унікальні активні ваучері
\item \$320,692 Об'єм обміну Пулів
\item 899 опубліковані звіти (\href{https://sarafu.network/reports}{архів історичних звітів})
\end{itemize}
Ці дані є історичними свідченнями, а не поточними даними використання CLC.

В більш широкому запропонованому проекті CLC досліджується, як сумісні мережі можуть:

\begin{enumerate}[leftmargin=*]
\item відкривати і цитувати шляхи через незалежні Пули;
\item координувати відповідальні мережні послуги, не переважаючи на місцеву відповідальність;
\item підтримка окремо фінансованих мандатів щодо ліквідності, гарантій, резервів або страхових полісів;
\item оцінити супи групи окремо від подачі, виконання та виплати ваучерів; і
\item використовувати запропоновану мережу та плату за послуги для фінансування прийнятих спільних послуг.
\end{enumerate}
Пропозиція охоплює \textbf{пропонований токен управління CLC} і \textbf{пропонований мережевий Пул CLC}. Жоден із них не є частиною поточного CLC App або Protocol v1.1.0.

За пропонованою моделлю учасники ліквідності та управління могли б діяти лише через окремо ухвалені програми з опублікованими вимогами до участі, ризиками, засобами контролю, правами передавання, умовами відновлення й правилами комісій. Звичайні депозити до Пулу не створюють автоматичного права на частку, повернення, виведення, винагороду або управління.

Об'єднаною метою є:

Підвищення відповідального обміну та виконання реальних зобов'язань, зберігаючи при цьому піклування, справедливість, місцеву відповідальність та стійкість.

Протифік захоплення і надійний вихід залишаються основними проектованими цілями. До запропонованих контрольних заходів відносяться затриманий час управління, багаторазові пороги затвердження, прозоре делегування, правила конфлікту, перегляд інцидентів та документований процес виведення і виведення. Ці контролі зменшували б ризики відбрання управління, але не могли б усунути втрати або гарантувати результати.


\section*{1. Протокол об'єднання зобов'язань (CPP): основна примітивна}
\addcontentsline{toc}{section}{1. Протокол об'єднання зобов'язань (CPP): основна примітивна}

\textbf{Модель розуму:} Пул зобов'язань - це регулювана угода для прийняття ваучерів або інших активів, публікації правил обміну, зберігання запасів та надання можливості для обмінів. Пізніше власники подають ваучерів своїм емітентам для реалізації в реальному світі. Обмін резервів і виконання емітентом - це окремі життєві цикли.

CPP координує цінність за допомогою чітко описаних зобов'язань. Модель обговорюється в \href{https://willruddick.substack.com/p/grassroots-economics-the-book-is}{Основна економіка: роздуми та практика}.

\subsection*{1.1 Що таке зобов'язання?}
\addcontentsline{toc}{subsection}{1.1 Що таке зобов'язання?}

Забезпечення - це обіцянка ідентифікованої сторони щодо майбутньої доставки, наприклад, їжі, транспорту, праці, зберігання або іншого законного товару, послуги, вигоду або продукції. А \textbf{ваучер} є ознакою або записом, представленим як ця зобов'язання за опублікованими умовами.

Токен-контракт записує цифрову механіку. Умови ваучера визначають емітент, Оферту, потужність, місце, час, обмеження, представлення, виконання, скарги та процес списання.

\subsection*{1.2 Що таке Пул зобов'язань?}
\addcontentsline{toc}{subsection}{1.2 Що таке Пул зобов'язань?}

Пул зобов'язань - це регулювана угода. Його може керувати індивідуальна особа, кооператив, громадська група, державне агентство, федерація, багатозначний оператор, оператор послуг або інша відповідальна структура.

До відповідних ролей і органів влади відносяться:

\begin{itemize}[leftmargin=*]
\item \textbf{Куратор Пулу:} публікує і керує правилами групи та будь-якими чітко передбаченими гарантіями;
\item \textbf{Власник Пулу:} володіє чинними повноваженнями власника \texttt{SwapPool};
\item \textbf{прокси-адміністратор:} може оновлювати проксидну реалізацію;
\item \textbf{контролерів залежності:} регулювати конфігуровані реєстри, котируючі котирування, обмежувачі або компоненти плати;
\item \textbf{знахідник маршруту або оператор:} може відкривати цитати або, в подальшому впровадженні, виконувати окремо затверджену маршрут; і
\item \textbf{гарант:} приймає на себе визначене зобов'язання лише за допомогою опублікованих, фінансованих умов.
\end{itemize}
CPP групи Функції Пулу розділяються на чотири поняття:

CPP об'єднує функції Пулу в чотири поняття:

\begin{itemize}[leftmargin=*]
\item \textbf{Курація:} допуск підтримуваних токенів або ваучерів.
\item \textbf{Оцінювання:} публікація методу, який використовується для визначення обмінного курсу або котирування.
\item \textbf{Обмеження:} застосування поточних лімітів балансу токенів у Пулі або інших окремо реалізованих засобів контролю.
\item \textbf{Обмін:} зберігання запасу, виконання обмінів, облік комісій і створення записів транзакцій.
\end{itemize}
Protocol v1.1.0 реалізує ці функції через \texttt{SwapPool} і необов'язкові залежності. Поточний \texttt{Limiter} обмежує баланс токена в Пулі; він не встановлює ковзних лімітів, лімітів на окремий обліковий запис або на всю мережу. Поточний \texttt{SwapRouter} обчислює котирування для кількох Пулів; він не виконує обміни.

\subsection*{1.3 Логіка поточного прямого обміну}
\addcontentsline{toc}{subsection}{1.3 Логіка поточного прямого обміну}

Сучасний прямий обмін Пул:

\begin{enumerate}[leftmargin=*]
\item перевіряє факультативний реєстр вхідних та вихідних токенів;
\item вимірює отримані дані;
\item отримує цитату з конфігурованого котиру або застосовує парату грубої одиниці;
\item перевіряє отриманий баланс токенів Пул проти факультативного лімітера;
\item розраховує плату групи та будь-які додаткові протокольні плати;
\item перевіряє наявний запас продукції;
\item переводить протокольну плату та вихід і облікує плату групи; і
\item видає обмінні події.
\end{enumerate}
Цитата є параметром транзакції, а не доказом потужності емітента, вартості погашення, справедливої вартості, грошової конвертації або гарантії.

\subsection*{1.4 Швидше використання}
\addcontentsline{toc}{subsection}{1.4 Швидше використання}

Пули зобов'язань може підтримувати обмін спільнотами, виробництво, взаємну допомогу, громадські програми та інші відповідальні структури. Документаційний кредитний продукт може використовувати ваучер як обезпеку або інструмент погашення, але для цього потрібні додаткові та конкретні умови для транзакції. Зазвичай відправка, депозит на базі, обмін на базі, подача погашення, виконання або виплата не є автоматично кредитом або погашенням.

CPP призначений для обмінної відповідальності та аудиторських записів транзакцій, а не для спекулятивних розрахунків. Записи на ланцюгу досі не доказують реалізації в реальному світі або соціального впливу.


\section*{2. Складання бухгалтерського обліку: від активів до довіри}
\addcontentsline{toc}{section}{2. Складання бухгалтерського обліку: від активів до довіри}

Традиційні фінанси починаються з \textbf{Активи - зобов'язання = власний капітал} CPP переформує це для економіки зобов'язань:

\begin{itemize}[leftmargin=*]
\item \textbf{Можливість прийняття:} кількість зобов'язань емітента, яку фонд або мережа все ще готові прийняти.
\item \textbf{Планувані зобов'язання:} Видав обіцянки, які залишаються не виконаними і виконуються іншими.
\item \textbf{Можливість виконання:} підтвердженою здатністю емітента виконувати ці обіцянки.
\end{itemize}
\subsection*{Ілюстративний зв'язок зобов'язання-здатність:}
\addcontentsline{toc}{subsection}{Ілюстративний зв'язок зобов'язання-здатність:}

Можливість прийняття - неплачувані зобов'язання = остаточна здатність прийняття

Ця концептуальна ідентичність може позначити обмеження ризиків групи: необмежене видання не означає необмежене прийняття. Це не ідентифікація балансу або заява про те, що кожен ваучер є юридично боргом або кредитом.

\textbf{Приклад:} Транспортер випускає 100100 поїздок  в ваучерів. Пул приймає максимум 40 поїздок за раз. Якщо 15 прийнятих поїздок залишаються незапланованими, Пункт має здатність прийняти до 25 додаткових поїздок за цією політикою. Вирахунок сам по собі не створює кредиту або виконання гарантії.



\section*{3. Діяльність виконання, списання та обміну}
\addcontentsline{toc}{section}{3. Діяльність виконання, списання та обміну}

Зміни фондового обміну, які стосуються активів державної позиції. Витримання ваучера відбувається, коли емітент виконує опубліковану зобов'язання. Вилучення записує виконані одиниці, щоб вони не могли бути представлені знову. Ці події не повинні бути зруйновані на одне використання поселення.''

У запропонованій системі вимірювання використовуються окремі записи для:

\begin{itemize}[leftmargin=*]
\item завершений обсяг обмінного фонду;
\item дійсні представлення виплати;
\item підтвердженої виконанню вимог емітента;
\item списання;
\item вистачивши кваліфіковані зобов'язання; і
\item вимірений запас Пулу.
\end{itemize}
Пропозиційна швидкість списання зобов'язань може порівнювати виконану вартість протягом періоду з середньостатистичною допустимою вартістю зобов'язань, але тільки тоді, коли обидва використовують один і той же метод розкритого оцінювання. Доставка токенів не є достатнім номінатором: вона може включати запас емітента, закінчені або недоступні одиниці, випробування та баланси, які не представляють вистачуючих зобов'язань третіх осіб.

Роздільний співвідношення супа-активності супу може порівняти завершену вартість супу супу з вимірюваним запасом супу. Він описує обмінну діяльність, а не виконання емітентом.

Ліквидність може покращити доступність запасів і полегшити обмін. Це не гарантує, що емітенти будуть виконувати свої зобов'язання, що вносичі будуть відновлювати активи, або що котирування групи представляє погашення або грошову вартість. Будь-які права постачальника ліквідності вимагають окремих опублікованих термінів.

Ви бачите \href{https://docs.cosmolocal.credit/uk/white-paper/appendix-a-math-box}{Додатку А} для визначення і \href{https://docs.cosmolocal.credit/uk/white-paper/appendix-c-kpi-definitions}{Додатку В} для запропонованої специфікації вимірювання.


\section*{4. Багаторазове забезпечення форвардного типу}
\addcontentsline{toc}{section}{4. Багаторазове забезпечення форвардного типу}

Роздільно задокументована кредитна установа може прийняти переносимі ваучери як заклад або як інструмент погашення. Якщо це сталося:

\begin{itemize}[leftmargin=*]
\item Прийняття фонду може зробити заклад більш відкритим;
\item додаткові варіанти законного представлення та виконання можуть підтримувати погашення;
\item обмеження, запаси, оцінка, ліквідності, продуктивність емітента та ризики від погашення зобов'язань залишалися б; і
\item жодного маршруту, виконання, повернення, вартості або відновлення не було б гарантовано.
\end{itemize}
\subsection*{4.1 Кредитний цикл виробника}
\addcontentsline{toc}{subsection}{4.1 Кредитний цикл виробника}

Майбуча послуга, що відповідає CLC-, могла б підтримувати кредитування виробників лише тоді, коли сторони встановлюють окремий механізм і прямо представляють операцію як відплачуваний аванс. У його умовах операції було б визначити кредитора та боржника і заявити:

\begin{itemize}[leftmargin=*]
\item суми авансу та погашення;
\item терміни погашення, відсотки, збори та дозволене використання закладних засобів;
\item передачу і будь-які зміни кредитора або власнику претензій;
\item як оцінюється і підтверджується виконання ваучера для бухгалтерського обліку погашення;
\item частине погашення, остатні суми, дефолт і розплата; і
\item доступні засоби захисту відповідно до чинного законодавства.
\end{itemize}
Ілюстративний потік:

\begin{enumerate}[leftmargin=*]
\item Виробник або постачальник послуг видає ваучер, що свідчить про зобов'язання щодо майбутньої продукції, наприклад, десять поїздок на таксі, 50 кілограмів кукурузи або десять робочих годин.
\item Куратор Пулу визнає цей ваучер і публікує метод обмінного курсу групи, обмеження, збори, правила запасу та будь-яку чітко передбачену охорону.
\item Кредитор або програма ліквідності окремо передбачає аванс за письмовими умовами операції та приймає ваучерів виробника як заклад або узгоджену інструмент погашення.
\item Власники можуть передати ваучери, обміняти їх через відповідні фонди або представити їх видавнику.
\item Потвердженої виконання емітентом гарантії задовольняється зобов'язання про ваучер. Він знижує окремий кредитний баланс лише в мірі, на вартістю і на основі доказів, зазначених у кредитних умовах.
\item Документи відрізняють виконання ваучерів від бухгалтерського обліку кредитів і визначають будь-які часткові погашення, остаточну суму, передачу, дефолт або виплату.
\end{enumerate}
Ця структура могла б дозволити узгодженому погашенню в натуральному вигляді, зберігаючи при цьому окремі записи про ваучер і кредитні зобов'язання. Вона не випливає з звичайного переказу, депозиту, обміну резерву, представлення погашення, виконання або списання.

\textbf{Приклад:} Продюсер отримує аванс в розмірі \$1,000 за умовами, які вказують, що підтвердженої виконання зазначених кукурузних ваучерів кредитується на суму повернення за допомогою зазначеного методу оцінки. Кредит застосовується кредитором тільки після отримання доказів, необхідних за цими умовами. Якщо кредитні умови не роблять такого зв'язку, придбання, передача, обмін, подача або виконання ваучера не впливає на аванс.


\section*{5. Від ізольованих Пулів до федеративної мережі}
\addcontentsline{toc}{section}{5. Від ізольованих Пулів до федеративної мережі}

Останній Protocol v1.1.0 підтримує безпосереднє виконання через один \texttt{SwapPool} і забезпечує тільки котировку \texttt{SwapRouter}. Він не виконує багатопропускні маршрути, HTLC, маршрути по зберіганню, партійное обчислення або розчищення між мережами.

У цій главі пропонується, як самостійно керувані Пули могли б координувати роботу, не відмовляючись від своїх власних правил вхвалення, оцінки, обмеження, зборів, запасу, уповноваження та управління.

\subsection*{5.1 Роздільні заходи обміну та виконання}
\addcontentsline{toc}{subsection}{5.1 Роздільні заходи обміну та виконання}

Федерація може поліпшити доступ до запасів, але вона не об'єднає ваучер і обмінять життєві цикли. Будь-яка реалізація оцінювала б ці події окремо:

\begin{enumerate}[leftmargin=*]
\item цитується маршрут;
\item один або більше обмінних акцій групи виконують і розраховують на ланцюжку;
\item власник представляє емітенту ваучерні одиниці;
\item емітент виконує зобов'язання; і
\item виконані одиниці розкідуються.
\end{enumerate}
Більше цитованих або виконаних маршрутів не свідчать про більшу реалізацію. У звітах було б зазначити коорту, період, активи, метод оцінки та часовий штамп, виключення, виправлення та докази поза ланцюгом, необхідні в Додатку С.

\textbf{Ілюстраційний маршрут:} Школа має ваучери для кукурузи, але їй потрібні транспортні ваучери. Сервіс маршруту ідентифікує сумісні запаси Пулів. Виконання було б успішним лише в тому випадку, якщо кожен окремо уповноважений хоп залишився в межах своїх квотів, обмежень, зборів, запасу та політики. Виникла супа не довела б, що будь-який з емітентів пізніше виконував свої зобов'язання щодо ваучера.

\subsection*{5.2 Пропановані послуги маршрутування та перерівноваження}
\addcontentsline{toc}{subsection}{5.2 Пропановані послуги маршрутування та перерівноваження}

Майбутні маршрутні послуги можуть підтримувати дві різні заходи.

\textbf{Виконання, яке було ініційовано учасником.} З урахуванням вхідних та вихідних активів, обсягу та обмежень користувача, послуга могла визначити шлях і підготувати виконання. У кожного хопу буде свій власний відповідальний Пул, цитат, авторизація, збори, обмеження, запас і квитанція. Атомічні партії, HTLC і гарантії є можливими варіантами виконання в майбутньому, а не поточним поведінкою протоколу.

\textbf{Оптимізація Пулу.} Куратори Пулів може опублікувати цілевизначення інвентарів, дозволені контрагенти, класи активів, обмеження відхилення котирувань і обмеження на період. Відповідальна служба могла б шукати сумісні цикли або ланцюги і виконувати тільки уповноважені наміри.

Ребалансування було б оптимальним. Пункт може дозволити маршрути учасників, відмовившись від ребалансування, або він може дозволити тільки вибрані активи, контрагенти та суми. Кожна виконана стяжка отримала б квитанцію, а будь-яка плата за обслуговування була б розкрита окремо від плати за фонд та протокол.

\subsubsection*{5.2.1 Конфедерація та оперативна здатність}

Незалежні розгортання могли б управляти власними реєстрами, інтерфейсами, маршрутними службами та профілями політики, вибираючи при цьому сумісні стандарти даних та отримання. Виконання міжпрофілів залишилося б залежним від розгортання.

Сумісний профіль:

\begin{itemize}[leftmargin=*]
\item ідентифікують свої коріння реєстру, операторів послуг, контролерів та відповідні терміни;
\item розкривати дозволені і відмовлені контрагенти, активи, адаптери та маршрути;
\item Використовують у всіх учасниках групи санкції, обмеження, збори та обмеження на запаси;
\item зберігати доказники від пропозиції до отримання на хоп; і
\item дозволяти іншим чином функціональним пулям залишати або вибирати інший реєстр без видалення балансів або зобов'язань емітента.
\end{itemize}
Сумішність може збільшити доступні шляхи обміну і зменшити залежність від одного реєстру або оператора. Це не робить мережу, CLC App, GEF або інший фонд відповідальним за виконання емітентом.

\subsection*{5.3 Пропонована модель частки мережі та сервісної комісії}
\addcontentsline{toc}{subsection}{5.3 Пропонована модель частки мережі та сервісної комісії}

Поточний Protocol v1.1.0 стягує Комісію Пулу та, якщо її налаштовано, додаткову Протокольну комісію за прямий Обмін через Пул. Ці поточні комісії залишаються окремими.

Майбутня програма могла б окремо отримувати:

\begin{enumerate}[leftmargin=*]
\item \textbf{частку мережі}, визначену як оголошена частка зібраних Комісій Пулів-учасників; і
\item \textbf{маршрутну або сервісну комісію}, що стягується за визначений майбутній сервіс.
\end{enumerate}
Пропонована частка мережі не є додатковим відсотком від повної суми обміну після врахування Комісії Пулу. Для Пулу \texttt{p}:

\[
\tau_p = f_p \cdot r_p
\]

де \texttt{f\_p} --- ставка Комісії Пулу, а \texttt{r\_p} --- пропонована частка цієї комісії, що спрямовується мережевій програмі.

За виміряний період:

\begin{itemize}[leftmargin=*]
\item \textbf{валові Комісії Пулів} --- сума фактично зібраних Комісій кожного Пулу;
\item \textbf{надходження від частки мережі} --- оголошені частки цих зібраних комісій;
\item \textbf{надходження від сервісних комісій} --- окремо зібрані маршрутні або сервісні комісії; і
\item \textbf{комісійні надходження програми} дорівнюють надходженням від частки мережі плюс надходження від сервісних комісій.
\end{itemize}
Жодна категорія не враховується двічі. Поточні Протокольні комісії не включаються, якщо окремо ухвалена політика законно не спрямовує фактичні надходження від Протокольної комісії до майбутньої програми.

Для сукупної оцінки припустімо:

\begin{itemize}[leftmargin=*]
\item \texttt{Q\_swap} --- вартість виконаних Обмінів через Пули для визначеної когорти й періоду; і
\item \texttt{$\tau$} --- ефективна ставка пропонованої частки мережі та окремо визначених сервісних комісій на цю виконану вартість Обмінів.
\end{itemize}
Тоді:

\[
F \approx \tau \cdot Q_{\mathrm{swap}}
\]

Це аналітична оцінка, а не обіцянка доходу. Кожне вхідне значення потребує визначених когорти, періоду, одиниці, часу оцінювання, виключень і політики виправлень.

\subsubsection*{5.3.1 Надходження, придатні для грошових виплат, і натуральні надходження}

Комісії можуть надходити у взаємозамінних активах, придатних для грошових виплат, або у ваучерах та інших натуральних активах. Натуральні надходження не можуть автоматично покривати грошові витрати чи вимоги щодо покриття. Будь-який обмін або конвертація потребуватиме відповідних повноважень, доступного запасу, розкритого маршруту, дотримання лімітів і фактичного виконання.

Нехай \texttt{$\chi$} позначає реалізовану частку комісійних надходжень, придатну для грошових виплат після застосування політик, невдалих конвертацій і прослизання. Тоді такі надходження становлять:

\[
F_{\mathrm{cash}} \approx \chi \cdot F
\]

Для аналізу бюджету й беззбитковості слід використовувати реалізоване значення \texttt{F\_cash}, а не валові нараховані комісії чи номінальну вартість натурального запасу. Майбутні програми мають окремо звітувати про валові комісії, натуральні надходження, спроби конвертації, реалізовані грошові надходження, витрати та залишки.

\subsection*{5.4 Пропановані програми ліквідності}
\addcontentsline{toc}{subsection}{5.4 Пропановані програми ліквідності}

Майбутні, окремо документовані програми щодо ліквідності можуть виділити активи на призначені Пулли або маршрутні послуги. Сучасні договори \texttt{SwapPool} не чекають акції Пул або автоматично не створюють права на погашення, виведення, винагороду, управління або прибуток.

Будь-яка програма опублікувала б:

\begin{itemize}[leftmargin=*]
\item відповідальний суб'єкт та учасник Куратори Пулів;
\item активів, які внесли внесок, і те, чи є передача погасним, витягуючим, дарованим або даруваним;
\item умови опіки та технічного контролю;
\item дозволені використання, обмеження, замки, ворота виведення та розподіл втрат;
\item виплата або права на стимули і чи може сума бути нульовою;
\item звітування, конфлікти, скарги та засоби захисту; і
\item міграція, припинення та обробка залишених активів та зобов'язань.
\end{itemize}
Суттєві ризики включають перелік, який важко обміняти або виконувати, низьку грошову допустимість, невиконання емітентом, невдачі договорів або постачальників, зміни управління та обмеження на вихід. Ограничення, резерви, квитки та підроздільні панелі можуть зменшити або розкрити деякі ризики; Вони не усуняють втрати.

Для аналітичної метрики ex-post, давайте:

\begin{itemize}[leftmargin=*]
\item \texttt{$\phi$} є реалізованою часткою прийомів програмних зборів, розподілених відповідно до прийнятих програмних умов; і
\item \texttt{K} - оцінена вартість активів, які охоплюються програмою за одним зазначеним методом.
\end{itemize}
Тоді:

\[
\mathrm{FeeFlow}_{LP} \approx \frac{\phi \cdot F}{K} = \frac{\phi \cdot \tau \cdot Q_{\mathrm{swap}}}{K}
\]

Ця метрика описує реалізований потік зборів за вимірений актив програми. Це не APY, прогноз, дивіденд або гарантований прибуток. У звітах буде зберігатися окремий обсяг свапів, представлення виплати, виконання емітентом, виплата, тривалість зберігання, втрати, виведення та отримання зборів.


\section*{6. Пропозиція ліквідності і управління на рівні мережі}
\addcontentsline{toc}{section}{6. Пропозиція ліквідності і управління на рівні мережі}

Досвід з історичного Sarafu Network спонукав дослідження двох більш широких питань дизайну:

\begin{enumerate}[leftmargin=*]
\item як самостійно управлені Пункти можуть координувати послуги з ліквідності та маршрутизації; і
\item як спільні послуги можуть залишатися відповідальними, оскільки участь зростає.
\end{enumerate}
Один із пропонованих дизайнів CLC передбачає \textbf{пропонований мережевий Пул CLC} як механізм клірингу на рівні мережі й координації бюджету та \textbf{пропонований токен управління CLC}. Жодного з них немає в поточному CLC App або Protocol v1.1.0. Інші розгортання можуть використовувати кооперативи, державні органи, федерації, мультипідписи, операторів сервісів або інші відповідальні структури без прийняття жодної з цих пропозицій.

\subsection*{6.1 Контроль і опціональність}
\addcontentsline{toc}{subsection}{6.1 Контроль і опціональність}

В даний час Protocol v1.1.0 може застосовувати криптовалюти балансу токенів Пул та перевірки запасів. Ці контролю обмежують вибрані контрактні дії; вони не запобігають погашенню результатів діяльності емітента, втраті вартості, втраті ключів, невдачі договору або будь-якій іншій втраті.

У подальшому розпорядженні можуть бути додані переломники, часові блоки, резерви, гарантії або страхування. Будь-яка охорона потребує ідентифікованої відповідальної сторони, охопленої події, фінансування, обмеження, виключення, тривалості, доказів, процесу подання претензій та застосуваних умов.

Пропановане управління повинно зберігати вузькі та переглядові повноваження реєстру та служб. Звичайні дії можуть включати повідомлення, затримку, публікацію підстав, апеляцію та виправлення. Надзвичайні дії повинні створювати реєстрацію інцидентів і закінчуватися чи отримувати перегляд відповідно до прийнятої політики.

Сумісний пошук маршрутів може зробити більш можливим обмін між Пулами. Виконання багатопропускних операцій потребує додаткового програмного забезпечення, авторизації, обмежень, обробки провалів і управління. Кількість коцінованих або виконаних обмінних операцій сама по собі не підтверджувала б виконання ваучера або соціальний вплив.


\section*{7. Пропанова CLC активів управління та управління}
\addcontentsline{toc}{section}{7. Пропанова CLC активів управління та управління}

У цій главі представлено один факультативний проект майбутнього управління та координації ліквідності. Пропозиційний токен управління CLC, stCLC, sCLC, сховище управління, страховий фонд, Водопад, мізерники, мандати, вікна з зборів і кредитів та програми зовнішнього ринку не є частиною поточного CLC App або Protocol v1.1.0. Кожен з них потребує реалізації, відповідального оператора, опублікованої політики та правил, перегляду безпеки та відповідного законодавчого затвердження.

CPP не вимагає цієї моделі токенів. Сумісні розгортання можуть скористатися кооперативами, державними агентствами, федераціями, багатокомпанії, службами або іншими відповідальними структурами.

\subsection*{7.1 Мета}
\addcontentsline{toc}{subsection}{7.1 Мета}

Якщо це буде прийнято, запропонований рівень управління може:

\begin{itemize}[leftmargin=*]
\item координує спільні реєстри та послуги маршрутів;
\item розподіляти затверджені мандати щодо ліквідності між незалежними управлінними пулями;
\item управляти факультативною, чітко визначеною, фінансованою програмою покриття;
\item підтримувати спільну інфраструктуру та спостереження;
\item визнавати внеслих внесок згідно з опублікованими правилами щодо допустимості і антигульня; і
\item зберегти перегляд, апеляцію, прозорість та надійний вихід, оскільки участь зростає.
\end{itemize}
\subsection*{7.2 Пропанована модель управління активами}
\addcontentsline{toc}{subsection}{7.2 Пропанована модель управління активами}

Проект містить три окремі запропоновані активи:

\begin{enumerate}[leftmargin=*]
\item \textbf{запропонований токен управління CLC:} Передається базовий управлінський актив відповідно до прийнятих правил видачі, зберігання, голосування та відповідності.
\item \textbf{stCLC:} непередачую квитанцію голосового опіку, виготовлену при запереченні запропонованого токену управління CLC. Він представляв би обмежену владою управління і міг би дозволити розкриття делегацій.
\item \textbf{sCLC:} актива-дозвол або стимул, виданий в рамках політики, що має епоху. Вона може підтримувати обмежений доступ до платної кредиту або затверджені ліквідні та маршрутні стимули, і може закінчитися або бути спалена наприкінці епохи.
\end{enumerate}
Ні stCLC, ні sCLC не були б власним капіталом, дивидендом, остаточною вимогою, обіцянкою повернення або об'єднаним ваучером. Прийнята політика може покласти нуль на еквізицію sCLC та доступ до платно-кредитного кредиту в будь-яку епоху.

Перед використанням були б опубліковані блоки управління, мінімальні тривалість, розхладження, делегація, обмеження концентрації та пороги критичних дій. Розблокування запропонованих токенів управління CLC не буде голосувати за цим проектом.

\subsubsection*{7.2.1 Ілюстративне постачання та розподіл}

Ілюстративна загальна пропозиція CLC податка токенів управління \textbf{500,000,000} Це дизайнерський параметр, а не поточний випуск, пропозиція, оцінка або обіцянка ліквідності.

Ілюстраційний розподіл:

\begin{itemize}[leftmargin=*]
\item \textbf{15\% --- Grassroots Economics Foundation:} на постійній основі, непередається і пов'язана з ідентифікованим GEF мультисигом відповідно до опублікованих правил підписання та ротації.
\item \textbf{15\%  основна команда та перші партнери:} під час встановленої політичної скали і лінійного періоду отримання 24 місяців, з загальним розкриттям правил голосування та передачі.
\item \textbf{30\%  приватні нагороди:} визнання ранніх учасників, що надають затверджену ліквідітність мережі, з встановленим полісом на 24 місяці.
\item \textbf{40\%  державний резерв :} проведено в управлінському сховішці, що обмежується часом для факультативних програм зовнішнього ринку та доступності.
\end{itemize}
У рамках ілюстраційного контролю державної ліквідності:

\begin{itemize}[leftmargin=*]
\item не більше 10\% загального постачання буде активним на зовнішніх місцях одночасно;
\item кожна розгортання закінчується після 90 днів, якщо вона не буде відновлена за допомогою прийнятого процесу;
\item позиції щодо ліквідності та будь-які позиційні токенів залишалися б під контролем розкритого управління; і
\item запропоновані токенів CLC управління, які використовуються в позиціях на ліквідності, не проголосували б, якщо б вони не були заблоковані відповідно до тих же правил, що і інші токенів на голосування.
\end{itemize}
Майбутні запуски можуть розпочинатися з розкритого мультисигу і переходу тільки після окремо затвердження визначних етапів, таких як незалежні аудити, моніторинг, інцидентні рублі, а також перевірені процедури паузу та розв'язання. Кожен прийнятий параметр і процес зміни будуть опубліковані окремо.

\subsubsection*{7.2.2 Стадії доступності та нагородження}

У майбутньому програмі нагородження можна було б використовувати етапові вікна внеску та опубліковану референтну вартість для прийому та бюджетування. Це посилання не було б обіцянкою ринкової ціни, оцінки, погашення або виходу.

Умови внеску вказують, чи є активи подаровані, надані, витягуються або погашуються; хто їх контролює; їх дозволене використання; замки; розподіл втрат; звітування; і умови виходу. Великі внески могли би зіткнутися з обмеженнями голосування або зменшенням ваги управління.

Будь-яка ліквідуальність на зовнішньому ринку вимагала б окремого рішення, яке визначило б місця проведення, відповідальних операторів, активи, ліпіки, терміни, конфлікти, контроль ризиків, звітування та припинення. Політика не передбачає ціни або не гарантує ціни.

\subsubsection*{7.2.3 Пропанована програма впливу}

Майбуча програма могла б виділити частину сховища управління вкладчикам, які розміщують активи в визначеному Пули зобов'язань і замірно покращують доступ до бірж і підтверджену виконання вимог емітента.

Прийнята політика допустимості:

\begin{enumerate}[leftmargin=*]
\item визначити затверджені Пули, активи, мінімальну тривалість і терміни закриття;
\item визначити продуктивний внесок за допомогою доказів здійсненої супи та окремо підтвердженої презентації, виконання та списання;
\item оцінити маржовий внесок, а не загальну вартість, зафіксовану самостійно;
\item виключають самопродажу та циркулярну діяльність миття;
\item застосовувати максимальні розміри на одну суб'єктність та зменшуючий прибуток; і
\item надавати вікно спорів, спорів та поправок до кінцевого розподілу.
\end{enumerate}
Будь-яка грант залишатиметься залежно від опублікованих умов програми, правил отримання, допустимості, відповідності та втрат.

\subsection*{7.3 Пропозиція голосування, стимули та доступ до зборів}
\addcontentsline{toc}{subsection}{7.3 Пропозиція голосування, стимули та доступ до зборів}

Згідно з цим проектом, власники stCLC могли голосувати за підходячі Пули, мандати маршрутних послуг, спільні бюджети або інші прийняті пропозиції. Система курированої оцінки може перетворити голоси на обмежені стимули sCLC для кваліфікованих операторів ліквідності або маршрутування.

Політика оцінки підтримувала б виконані обмінні позиції і виконання емітентом окремо. Це не було б нагородою цитованих, але не виконаних маршрутів, не вирішених презентацій, самостійного продажу або загальної вартості, зафіксованої без доказів корисної діяльності.

Прийнятий процес управління також міг би опублікувати епохійний бюджет з зборів і кредитів. Кваліфіковані власники stCLC можуть отримати sCLC, що дозволяє обмежені, обмежені часові обмінники з призначених платіжних фондів на виділені активи або програми. Це буде політичний доступ до доступних запасів, а не пасивний дохід або володіння платіжним фондом.

\subsubsection*{7.3.1 Виконання зборів за послуги та вибір профілю}

У майбутньому профіль реєстру може вимагати, щоб учасники Пулів або маршрутних послуг використовували сумісний адаптер або гачок. Профіль може відмовитися від відкриття, маршрутизації або підтримки програми, якщо потрібна квитанція не підтверджує збору плати.

Такі назви, як \textbf{PoolFactory} або \textbf{FeeHook}, описують можливі майбутні модулі; це не контракти Protocol v1.1.0. Кожна реалізація має розкрити код, адреси, контролерів, отримувачів, ставки, допустимі транзакції, обробку збоїв і статус аудиту.

Виконання повинно бути конкретним. Пункт, виключений одним профілем, може продовжувати функціонувати на місцевому рівні або приєднатися до іншого сумісного реєстру, якщо його договори та залежності залишаються функціональними. Клієнти повинні розкрити доступні профілі і показувати відповідні збори перед оформленням дозволу.

\subsubsection*{7.3.2 sCLC бюджет і факультативні зовнішні контрольні системи плаву}

Після прийняття фінансування бюджетів з більш високим пріоритетом, майбутній процес міг би опублікувати епохи такси-кредит бюджету \texttt{F\_epoch}, включаючи нуль. Політика може призначити ліміт учасника пропорційно голосовій можливісті stCLC:

\[
\mathrm{limit}_{\mathrm{user,epoch}} = F_{\mathrm{epoch}} \times \frac{\mathrm{stCLC}_{\mathrm{user}}}{\mathrm{stCLC}_{\mathrm{total}}}
\]

Кожне вікно залишалося б підлягати перелікам дозволів, ліпенням, запасам, допустимості, правилам інцидентів і чинному законодавству.

Віддельна політика могла б дозволити обмежену, временно обґрунтовану придбання запропонованого токену управління CLC лише для зменшення вимірюваної зовнішньої поверхні атаки на управління. Це вимагатиме:

\begin{itemize}[leftmargin=*]
\item опублікований триггер, заснований на зовнішньому плаву протягом зазначеного періоду;
\item максимальна частка реалізованих, допустимих зборів і обсягу проведення;
\item виведення за часом без точного оголошення часу;
\item автоматичне припинення, коли не виконано прийнятої охоплення або діючих порогів;
\item відставлення або розміщення придбаних токенів на розкритому безголосовому рахунку; і
\item ясна заява про те, що програма не спрямована на ціну або не гарантує її і не впливає на виконання вимог емітента.
\end{itemize}
Політика може залишатися нездатною на невизначений час.

\subsubsection*{7.3.3 Portfolio Пули та атестації}

А \textbf{Портфель} було б звичайним Пул зобов'язань, що було організовано навколо місії або сфери служби. Його Куратор Пулу може бути індивідуальним, кооперативним, громадською групою, державним агентством, федерацією, мультисиг, оператором послуг або іншою відповідальною структурою.

Підтримка може відбутися через:

\begin{enumerate}[leftmargin=*]
\item прямі внески за публікуваними умовами фонду;
\item тимчасово обмежені мандати щодо ліквідності, затверджені за допомогою прийнятого процесу управління; або
\item sCLC спрямований доступ до обмеженого бюджету після водопаду, коли цей механізм включений.
\end{enumerate}
Портфельні фонди залишалися б незалежними і могли б використовувати спільні або незалежні реєстри.

Сертифікати могли б надати інформацію про прийом або ставлення до ризику, але не створювали б права на збори, прибуток або остаточні активи. Розгортання може використовувати:

\begin{itemize}[leftmargin=*]
\item \textbf{сертифікати, пов'язані з реєстрацією} від ідентифікованого перевірителя, який застосовує опублікований метод; або
\item \textbf{ваучерів перевіреного обслуговування} що являє собою погашується зобов'язання щодо здійснення зазначеного аудиту або оцінки.
\end{itemize}
Пункт розкрив би, як атестація впливає на прийом, методи обмінного курсу, обмеження або допустимість і як розглядаються помилки, закінчення терміну дії, конфлікти та апеляції.

\subsection*{7.4 Пропозиція водоспаду і бюджети}
\addcontentsline{toc}{subsection}{7.4 Пропозиція водоспаду і бюджети}

Пропозиційний водопад виділить лише реалізовані, допустимі квитки в рамках прийнятого процесу управління. Це розрізняє:

\begin{itemize}[leftmargin=*]
\item \textbf{грошово допустимі активи (\texttt{E\_cash}):} перераховані фондові активи, які можуть бути використані або перетворені в рамках політики для грошових зобов'язань; і
\item \textbf{актив у натуральному вигляді (\texttt{E\_kind}):} ваучерів або інших активів, які можуть підтримувати обмін у мережі, але не включаються в покриття або операційні зобов'язання у грошовій валюті, якщо вони не були фактично конвертовані.
\end{itemize}
Політика конверсії визначала б авторизованих операторів, активи, місця проведення, джерела цін, вікна часу, межі сльоту, обмеження, конфлікти, записи та процедури інцидентів. Конверсія казначейства не визначала б зобов'язання емітента щодо ваучера або обмінного курсу групи.

Ілюстраційний порядок пріоритетів:

\begin{enumerate}[leftmargin=*]
\item \textbf{Мета фінансованого покриття:} побудувати будь-який прийнятий резерв або запропонований страховий фонд до своєї розкритої мети для визначених покритих подій.
\item \textbf{Основні операції:} фінансувати обмежений, розкритий бюджет на юридичну роботу, освіту, комунікацію, інфраструктуру, аудити та обсервативність.
\item \textbf{Мандати щодо ліквідності:} розподіляти затверджені активи для зазначених цілей обслуговування Пулу або маршруту в рамках перегляду ліпіків та заходу сонця.
\item \textbf{Факультативне управління зовнішнім плаванням:} Фонд тільки тоді, коли його опубліковані триггерні та пріоритетні пороги задоволені; в іншому випадку виділити нуль.
\item \textbf{Пропанований бюджет такси-кредиту:} розподіляти будь-який затверджений залишок на призначені Пункти з урахуванням зборів і публікувати \texttt{F\_epoch}, який може бути нулевим.
\end{enumerate}
Бюджетні рішення можуть розглядати підтвердженої виконання, покритої експозиції, допустимих резервів, обмеження використання, виконаних маршрутів, концентрації, інцидентів та виконання гарантом. Кожна метрика буде дотримуватися доказів і правил кохорти в Додатку С.

Можливим потоком учасників було б:

\begin{enumerate}[leftmargin=*]
\item вносники передають кваліфіковані активи за окремо прийнятими умовами нагородження або програми;
\item допустимі учасники отримують і, якщо це дозволяється, зачиняють запропоновані токенів управління CLC для отримання stCLC;
\item реалізовані допустимі квитки або квитки за плату за обслуговування входять в водопад;
\item фонди водоспаду прийняли пріоритети відповідно; і
\item тільки активизована політика після водного попадку випускає sCLC або відкриває вікно кредиту з платами з обмеженими межами.
\end{enumerate}
Жоден крок не створює права власності, погашення, виведення, покриття, винагороди або управління за межами тих, які прямо зазначені в застосованих умовах.


\section*{8. Технічний обсяг та зростання}
\addcontentsline{toc}{section}{8. Технічний обсяг та зростання}

У цій главі описані обов'язкові або запропоновані роботи. Це не перелік характеристик, які гарантуються в поточному CLC App або Protocol v1.1.0.

Можливі сфери роботи включають:

\begin{itemize}[leftmargin=*]
\item Виконання маршрутування між сумісними пулями, з реєстрацією, котируванням, лімітом, платою та відкриттям запасів;
\item адаптери хранения або HTLC для перекрестного виконання доменів, де атомне розрахунки недоступні;
\item інтерфейси та інструменти політики для невеликих або персональних Пулів;
\item аудиторські реєстри ваучерів, Пулів, методів обмінного курсу, лимитов, контролерів та зборів;
\item інтеграції спеціалізованих платіжних постачальників, потоки виведення коштів і перевірки права на участь;
\item зв'язки з функціональними активами з зовнішньою місцею ліквідності для ребалансування та ліквідності платежів; і
\item корекція казначейства, яка охоплюється політикою, для прийнятого покриття, операційних витрат або мандатів на ліквідітність.
\end{itemize}
Ціни на зовнішньому ринку не визначають, що зобов'язаний емітент за умови ваучера. Пункт може використовувати захищену зовнішню посиланню для функціонального активу, але його опублікований метод обмінного курсу, обмеження, збори та запас будуть регулювати його котирування.

\subsection*{8.1 Пропановані норми маршрутного обслуговування і SDK}
\addcontentsline{toc}{subsection}{8.1 Пропановані норми маршрутного обслуговування і SDK}

\textbf{Відкриття.} Пропанований маршрутний сервіс буде запитувати ідентифіковані реєстри для введення активів, методи обмінного курсу, обмеження, збори, запас, інциденти та інформацію про контролера. Заховувані записи включали б межі свіжості та ідентифікатори джерела.

\textbf{Профілі мережі.} Клієнт може підтримувати більше ніж один корінний реєстр або профіль політики. Вона повідомляла б учасника про те, який профіль, контрагенти, адаптери, обмеження та відповідальні оператори послуг використовується цитатою. Транспрофільний маршрут повинен задовольняти всі відповідні умови хопу.

\textbf{Політика шляху.} Відповідний оператор може виключити небезпечні залежності або контрагенти і застосувати обмеження рівня маршруту, вимоги свіжості та критерії охорони здоров'я. Ці сигнали підтримали б рішення; вони не гарантують виконання або захист від втрат.

\textbf{Плата і обмеження.} Цю цитату об'єднали б платіжні збори групи, будь-які додаткові поточні протокольні збори та будь-які окремо запропоновані платіжні або послуги маршрутування. Виконання відмовляло б від скасованих цитат або порушеного кордону.

\textbf{Атомність і відновлення.} Виконання в кількох штатах було б атомним, якщо це було б можливо. Якщо вона використовувала HTLC або гарантії, служба розкривала відклади часу, маршрути аборту, відповідальних контролерів, процедури інцидентів та рештальні ризики.

\textbf{Пропанова орієнтування і відновлення рівноваги.} Опти-ін-сервіс може збирати наміри відновлення балансу і шукати сумісні цикли або ланцюги. Це було б:

\begin{enumerate}[leftmargin=*]
\item опублікувати машиночитану квитанцію з ідентифікацією виконаних циклів, активів, сумів, часових штампів оцінки та зборів;
\item виконувати прийняті максимальні обмеження на період та політику контрагентів;
\item відмовлятися від діяльності, яка порушує будь-які авторизації, обмеження або доступний інвентар Пул; і
\item зберігати детерміністичні дані та доходи для перегляду та розв'язання спорів.
\end{enumerate}
\textbf{Вимоги до SDK.} SDK для виконаних маршрутів забезпечить детерміністське відображення пропозиції до отримання, перевірки невимірності на хоп, зрозумілі коди провалу і аудиторські журнали. У поточному Protocol v1.1.0 \texttt{SwapRouter} представлені лише цитати; вона не виконує ці запропоновані маршрути.

\subsubsection*{8.1.1 Специфікація мінімальної сумісності конфедерації}

Екосистема Пул, яка шукає кроспрофільне маршрутизацію, опублікувала б машинно читуючу інформацію для:

\begin{enumerate}[leftmargin=*]
\item \textbf{Коріння реєстру:} Ідентифікатори активів, Пулів, методів обмінного курсу, обмежень та політики зборів, або одна корень, яка їх детерміністично вирішує.
\item \textbf{Квітки:} Профіль, активи, котируються і виходять, суми, джерело котирування і часовий штамп, ліміт-снайп, збори, результат запасу та результат виконання за кожен скок.
\item \textbf{Операційні сигнали:} свежими даними про запаси, обмеження використання, інциденти та будь-яку окремо підтверджену виконання або фінансовану охорону.
\item \textbf{Політичні обмеження:} Дозволені або відмовлені контрагенти, класи активів, адаптери та будь-які вимоги до страхованості.
\item \textbf{Коди провалу:} детерміністські пояснення відхилення, закінчення терміну дії, обмеження, запасу, політики, залежності або інцидентних невдач.
\end{enumerate}
Профіль може додати охоплення, відповідність, арбітраж або інші послуги, не роблячи їх вимогами для базової сумісності CPP. Кожна додаткова служба визначала б свою відповідальну сторону, повноваження, сферу діяльності та терміни.

\subsection*{8.2 Ліцензування, перевірка та вихід}
\addcontentsline{toc}{subsection}{8.2 Ліцензування, перевірка та вихід}

Protocol v1.1.0 контракти EVM-сумісні. Договори в каталозі \texttt{src} реєстрації Протоколу публікуються під номером AGPL-3.0, за винятком ідентифікованих не модифікованих компонентів третіх сторін, які зберігають свої власні умови. Опубліковані джерела, АБІ та інструкції щодо розгортання підтримують незалежний огляд, але самі по собі не доказують аудиту, безпечне розгортання або відповідність законодавству.

Кожен розгортання буде окремо розкривати свою версію коду, походження конструкції, адреси, повноваження контролера та оновлення, статус аудиту, дзеркали реєстру, а також будь-які захисті від блокування часу або паузу.

Пропозиція \textbf{набор з виліками} може включати:

\begin{enumerate}[leftmargin=*]
\item скрипти детерміністичного розгортання;
\item реєстр і інструменти експорту;
\item документований процес переведення маршрутних послуг, SDK та інтерфейсів на новий корінь реєстру;
\item перелік Куратор Пулу для безпечного виходу з спільного реєстру; і
\item перелік перемігрантів для неплачуваних ваучерів, включаючи повідомлення про емітент, терміни подачі та виконання, продовження доступу до записів та засоби захисту.
\end{enumerate}
Сумісні виліки можуть поліпшити стійкість, коли громади, кооперативи, державні органи, федерації, мультисіги або оператори послуг потребують іншого управління. Фактична безперервність все ще залежить від власності на контракт, ключів, залежності, інтерфейсів, інфраструктури, юридичних зобов'язань та послуг третіх осіб.


\section*{9. Предлагана економіка для програм ліквідності}
\addcontentsline{toc}{section}{9. Предлагана економіка для програм ліквідності}

Ця глава описує майбутню модель, прийняту окремо. Це не актуальна функція додатку, пропозиція, обіцянка повернення або право, створене шляхом вкладу в \texttt{SwapPool}.

\subsection*{9.1 Пропановані джерела доходів}
\addcontentsline{toc}{subsection}{9.1 Пропановані джерела доходів}

Майбутній бюджет мережі може отримати:

\begin{enumerate}[leftmargin=*]
\item розкритий \textbf{мережний рейк} прийнято як частку збірних зборів учасних Пулів;
\item окремі збори за маршрутизацію або послуги від реалізованих спільних послуг; і
\item інші прямо прийняті, отримані доходи.
\end{enumerate}
Об'єктний об'єкт об'єкта, що зберігається в об'єктах, не є доходом мережі. Останній Protocol v1.1.0 використовує іншу модель: факультативна плата за протокол додається до плати групи і відправляється безпосередньо до її конфігурованого одержувача.

\subsection*{9.2 Приклад розрахунку частки мережі}
\addcontentsline{toc}{subsection}{9.2 Приклад розрахунку частки мережі}

Якщо учасник-поєдинник взимає плату поєдинку 2.00\%, а прийнята мережна плата отримує 20\% цієї плати поєдинку, запропонована ефективна ставка поєдинку поєдинку на маршрутовану вартість становить:

\texttt{rake\_rate = 2.00\% $\times$ 20\% = 0.40\% = 40 bps}

Якщо 25\% отриманих активів за раке і плату за послуги є допустими і конвертируються для зазначеного використання в грошовій валюті після витрат, доля грошових коштів, що використовуються для раке за 40 кВт, становить приблизно 10 кВт.

Пропановані доходи від мережі складають:

\texttt{network\_rake\_received + routing\_or\_service\_fees\_received}

Не додавайте до мережного рейку брутні збори за Пул: рейк є переведенням з цих зборів і в іншому випадку буде розрахуватися двічі.

\subsection*{9.3 Права на програми ліквідності}
\addcontentsline{toc}{subsection}{9.3 Права на програми ліквідності}

Подробно прийнята програма може фінансувати запаси Пулів, маршрутні послуги, моніторинг або інші мандати. Його умови повинні розкривати:

\begin{itemize}[leftmargin=*]
\item чи є передачу подарунком, нагородженням, кредитом, відновлюваним внеском або покупкою;
\item опіку і контроль;
\item правила виведення, погашення, втрати та пріоритетів;
\item виплати та отримання винагороди;
\item права управління;
\item методи оцінки та звітності; і
\item призупинення, припинення і виправлення.
\end{itemize}
У поточному \texttt{SwapPool} не створюється токен Пул-share або автоматичний право на внесок. Будь-яка метрика ex-post повинна ґрунтуватися на реалізованих доходах і втратах, не повинна бути представлена як обіцяна вигода, і може бути нульовою або негативною.


\section*{10. Комплексна система ризиків}
\addcontentsline{toc}{section}{10. Комплексна система ризиків}

Ця запропонована рамка розділяє десять категорій ризиків. Для кожної категорії він визначає можливі показники, контрольні заходи, стрес-тести та показательний апетит до ризику. Це рекомендації про конструкцію, а не твердження про те, що кожен контроль використовується, ефективний або достатній. Ограничення, резерви, гарантії, моніторинг, охоплення та управління не можуть усунути втрати.

\subsection*{10.1 Протокол і ризик розумних контрактів}
\addcontentsline{toc}{subsection}{10.1 Протокол і ризик розумних контрактів}

\begin{itemize}[leftmargin=*]
\item \textbf{Погрози:} контрактні помилки, помилки модернізації, невдачі в залежності та неправильна конфігурація обмежень або зборів.
\item \textbf{Показники:} висновки аудиту, незрозумілі переміщення запасів, незмінні збої та незвичайні повернення транзакцій.
\item \textbf{Можливості контролю:} незалежні аудиториї; мінімальні привилегировані ролі; розкриті прокси і контролерів залежності; оновлення, встановлені вчасно; моніторинг на ланцюжку; інцидентні паузи з опублікованими повноваженнями та критеріями; і випробуваний міграційний шлях.
\item \textbf{Стрессові аналізи:} недоступні котируючі або обмеження залежності, дефіцити запасу, призупинення контрактів і вибух трафіку.
\item \textbf{Апетит до ризиків:} Низький обсяг перед скасуванням неплачуваних зобов'язань або обсягу обмінних операцій.
\end{itemize}
\subsection*{10.2 Економічні та ринкові ризики}
\addcontentsline{toc}{subsection}{10.2 Економічні та ринкові ризики}

\begin{itemize}[leftmargin=*]
\item \textbf{Погрози:} тонкий запас, односторонні потоки, швидкі виведення і маніпульовані або застарілі посилання на ціни.
\item \textbf{Показчики:} Високе використання балансу токенів, розширення розбіжностей котирувань, часті відхилення лимітів і концентрований запас.
\item \textbf{Можливості контролю:} поточні криптовалютні криптовалюти балансу токенів групи; запропоновані обмеження для розрахунку або рахунку; окремо прийняті резерви; захищені посилання на ціни; виключення маршрутів; а також стяги або обмеження інцидентів, що обмежені часом.
\item \textbf{Стрессові аналізи:} великі переміщення референтних цін, зростання представництва, запити на виведення і відключення джерел даних.
\item \textbf{Апетит до ризиків:} Модерувати лише за опублікованими параметрами та фінансованою здатністю нести втрати.
\end{itemize}
\subsection*{10.3 Ризик видавниця та ваучера}
\addcontentsline{toc}{subsection}{10.3 Ризик видавниця та ваучера}

\begin{itemize}[leftmargin=*]
\item \textbf{Погрози:} випуск за межами спроможності виконання, погашення зобов'язань випускника, помилкові умови та слабо визначені вікна доступності або презентації.
\item \textbf{Показчики:} зниження ставок підтвердженої виконання, застарівання нерешених ваучерів, нерешенних скарг та концентрована експозиція одному емітенту.
\item \textbf{Можливості контролю:} належну дбайність емітента; чіткий ваучер і умови пропозиції; обмеження на видачу або допуску; облигації або гарантії, що фінансуються самостійно; звітування з коортів; і відповідального перегляду реєстру.
\item \textbf{Стрессові аналізи:} банкрутство емітента, регіональні шоки виробництва, фальсифіковані вимоги та тривалі затримки у виконанні.
\item \textbf{Апетит до ризиків:} знижується з збільшенням концентрації емітента або пропозиції.
\end{itemize}
\subsection*{10.4 Ризик представлення і виконання виплати}
\addcontentsline{toc}{subsection}{10.4 Ризик представлення і виконання виплати}

\begin{itemize}[leftmargin=*]
\item \textbf{Погрози:} недійсне або подвійне представлення, недостатній потенціал емітента, запаси, логістичні невдачі та відсутність записів про списання.
\item \textbf{Показчики:} затримка виконання запитів, невдалені або суперечливі представлення, запаси, запізники квитків і виконані одиниці, які не мають списання.
\item \textbf{Можливості контролю:} опубліковані процедури представлення та виконання; розкриття потужностей; багаторічні місця виконання, якщо це законно; стандарти доказів; шляхи подачі скарг та виправлення; і записи, що розділяють подачу, виконання і списання.
\item \textbf{Стрессові аналізи:} два-четири рази більший обсяг подачі, переваги установки, збитки постачальників та спроби повторного використання.
\item \textbf{Апетит до ризиків:} Низький рівень, коли є пов'язані з необхідними товарами, вразливими учасниками або довгими вікнами виконання.
\end{itemize}
\subsection*{10.5 Ризик управління}
\addcontentsline{toc}{subsection}{10.5 Ризик управління}

\begin{itemize}[leftmargin=*]
\item \textbf{Погрози:} Куратор Пулу або контролер захоплення, швидкі зміни параметрів, конфлікти інтересів, приховані технічні повноваження, і слабкі апеляції.
\item \textbf{Показчики:} концентрирована влада, часті надзвичайні дії, необгрунтовані зміни політики та повторні переваги.
\item \textbf{Можливості контролю:} розкриття ролі та повноважень; пропорційні порожники затвердження; часові періоди; конфлікти та правила відмовлення; публічні записи про зміни; апеляції; автоматичні сонячні заходи з надзвичайною енергією; і витворюваність.
\item \textbf{Стрессові аналізи:} Пропозиції протистояння, втрата підписників, спроби підкупу і захоплення за допомогою накопичення або делегованого контролю.
\item \textbf{Апетит до ризиків:} низький для дій, що впливають на методи оцінки, виведення, коріння реєстру, охоплення або надзвичайні повноваження.
\end{itemize}
Для майбутнього розгортання запропонованого токену управління CLC тест захоплення включав би накопичення, послідовну спробами перенаправлення бюджетів, послаблення стандартів реєстру, затвердження мандатів пов'язаних сторін або вичерпнення фінансованого покриття. До можливих гарантій відносяться блоки управління, більш високі пороги для критичних дій, затримка виконання, прозоре моніторинг делегації та концентрації, процес інцидентів та надійний порядок виходу. Жодна з них не представлена як розгорнута лише тим, що з'явилася тут.

\subsection*{10.6 Правовий ризик та ризик відповідності}
\addcontentsline{toc}{subsection}{10.6 Правовий ризик та ризик відповідності}

\begin{itemize}[leftmargin=*]
\item \textbf{Погрози:} ваучер, послуга, просування або управлінський актив, який отримує несподіване регуляторне ставлення; провалу захисту споживачів; відмивання грошей або вплив санкцій; і транскордонні обмеження.
\item \textbf{Показчики:} Флаги юрисдикції, запити регулятора, скарги, співвідношення обмежених сторін і розбіжності між рекламним і фактичним поведінкою.
\item \textbf{Можливості контролю:} перегляд за класом та юрисдикцією; точні розкриття; геоосновані інтерфейси; пропорційні перевірки допустимості або атестації; контролю за підвищення кваліфікації; документи про авторитет і прийняття; і чіткі відповідальні сторони.
\item \textbf{Стрессові аналізи:} обмеження юрисдикції, припинення постачальника, обов'язкова перекласифікація та наказ про призупинення функції або класу активів.
\item \textbf{Апетит до ризиків:} низький; обмежити або зупинити непідтримувану діяльність.
\end{itemize}
\subsubsection*{10.6.1 Правове позиціонування та обробка запропонованих токенів}

\begin{enumerate}[leftmargin=*]
\item \textbf{Перевірована інфраструктура:} Контракти Protocol v1.1.0 є несумісними з EVM- і їх джерело публікується на підставі ліцензій та винятків третіх осіб, визначених у реєстрі протоколу. Публікація дає можливість перегляду, але сама по собі не доводить аудиту, безпечне впровадження або відповідність законодавству. Кожен розгортання повинен розкрити свою версію коду, походження конструкції, адреси, повноваження контролера та статус аудиту.
\item \textbf{Пропанована позиція знака:} у цьому проекті запропонований токен управління CLC координує управління та доступ до політики. Вона не створювала б дивидендів, розподілу прибутку, залишкових прав або гарантії вартості або ліквідності.
\item \textbf{Порядкові запропоновані активи:} окремо реалізована політика може дозволити заблокувати запропонований токен управління CLC на монету stCLC і може дозволити визначити епоху sCLC. У прийнятих термінах необхідно було б точно визначити їх права, обмеження, термін дії, передачу та лікування.
\item \textbf{Повідомлення:} матеріали для будь-якого запропонованого розподілу CLC-токенів управління, stCLC або sCLC не повинні обіцяти прибуток, оцінку, пасивний дохід або гарантований доступ.
\item \textbf{Контроль юрисдикції:} розгортання може вимагати геофонізації інтерфейсу, атестацій для обмежених класів, обмежень просування, перегляду конкретних активів або контролю, які зупиняють запропоновану функцію.
\item \textbf{За повідомленням учасника:} Прийняті умови пояснюють, коли доступ може бути зменшений або відключений з юридичних, оперативних або ризикових причин, а також чи застосовується будь-яка компенсація або правозахист.
\end{enumerate}
\subsection*{10.7 Ризик маршрутування та перетину доменів}
\addcontentsline{toc}{subsection}{10.7 Ризик маршрутування та перетину доменів}

\begin{itemize}[leftmargin=*]
\item \textbf{Погрози:} Частинне виконання, застрянуті перебіги, експлуатації мостових або гарантійних коштів, застарілі котирування, інфляція маршруту і перехід до оголошених змін вартості.
\item \textbf{Показчики:} строк закінчення терміну дії маршруту, запізників на депозитах, відмінностей від котирування до виконання, повторних непотрібних переходів і інцидентів на мостах.
\item \textbf{Можливості контролю:} атомне виконання, якщо це доступно; консервативні HTLC або гарантійні часові дії; шлях та політику контрагентів; орієнтовані на рівні маршруту; детерміністське відображення котирування до отримання; і відповідальні оператори послуг.
\item \textbf{Стрессові аналізи:} Пункт-стоп, реорганізація ланцюга, відключення залежності, і один невдалий скок на запропонованому маршруті.
\item \textbf{Апетит до ризиків:} від низького до середнього рівня тільки для виявлених і контролюваних залежностей.
\end{itemize}
\subsection*{10.8 Ризик зберігання і управління ключами}
\addcontentsline{toc}{subsection}{10.8 Ризик зберігання і управління ключами}

\begin{itemize}[leftmargin=*]
\item \textbf{Погрози:} втрачені або порушені ключі, спірність підписника і неясна влада відновлення або адміністратора.
\item \textbf{Показчики:} аномальні підписи, зміни контролера, невдалені ротації та незвичайні виведення.
\item \textbf{Можливості контролю:} багатозначне або порожневе затвердження; ключі, що підтримуються обладнанням; відділення ролі; ротація підписника; запаси державних контролерів; контрольовані обмеження виведення; і перевірені процедури відновлення.
\item \textbf{Апетит до ризиків:} Низко.
\end{itemize}
\subsection*{10.9 Репутація та соціальні ризики}
\addcontentsline{toc}{subsection}{10.9 Репутація та соціальні ризики}

\begin{itemize}[leftmargin=*]
\item \textbf{Погрози:} помилкові заяви, шкідливі стимули, недоступні скарги, поганий досвід виконання, порушення конфіденційності та захист, який сприяє інсайдерів.
\item \textbf{Показчики:} скарги по кохортам, нерешенні суперечки, тенденції продуктивності емітентів, концентрація прибутку або втрати, а також відгуки спільноти.
\item \textbf{Можливості контролю:} розкриття на простой мові; маршрути оплаки та виправлення; доступні докази; прозоре звітування; перегляд інцидентів; і пропорційними санкціями за неправду.
\item \textbf{Апетит до ризиків:} низький рівень, з особливою увагою до постраждалих громад та вразливих учасників.
\end{itemize}
\subsection*{10.10 Ризик концентрації та фрагментації}
\addcontentsline{toc}{subsection}{10.10 Ризик концентрації та фрагментації}

\begin{itemize}[leftmargin=*]
\item \textbf{Погрози:} залежність від невеликої кількості емітентів, пулів, контролерів, постачальників, мереж або несумісних вилок.
\item \textbf{Показчики:} заходи концентрування державою-емітентом, фондом, запасом, контролером або постачальником послуг; провалу маршруту між кластерами; і критичні окремі залежності.
\item \textbf{Можливості контролю:} опубліковані межі концентрації; численні оператори, що відповідають; сумісні стандарти; незалежні реєстри; перевірені процедури виходу; і безпечна оперативна взаємодія.
\item \textbf{Стрессові аналізи:} втрата найбільшого емітента, групи, оператора, постачальника або коріння реєстру.
\item \textbf{Апетит до ризиків:} специфічні для розгортання та розкриті, з більш жорсткими обмеженнями для основних послуг або незамінними залежностями.
\end{itemize}

\section*{11. Механізми управління}
\addcontentsline{toc}{section}{11. Механізми управління}

Ця глава пропонує шаблон управління. Це не означає, що поточний CLC App використовує голосування за токеном управління, часові блоки, спільне страхування, процес вимог або будь-який контроль, описанний нижче. Кожна розгортання повинна визначити своїх дійсних приймачів рішень, органів влади, контрактів, процесів та політики.

\begin{itemize}[leftmargin=*]
\item \textbf{Конституційні цінності:} турбота про людей, турбота про навколишнє середовище, справедливість, взаємність, недомінування і стійкість.
\item \textbf{Типи пропозицій:} зміни плати, ліміту та індексу; мандати щодо ліквідності; Перелік Пулів та виведення; рішення про факультативне охоплення; і параметрові огради.
\item \textbf{Процес відповідальності:} Прийняття $\to$ оцінка $\to$ перегляд ризиків $\to$ затвердження $\to$ час, коли це необхідно $\to$ виконання. Захвалення може прийти від стюардів, кооперативів, державних органів, федерацій, мультисиг, голосування на ланцюзі або іншої розкритої та відповідальної структури.
\item \textbf{Прагові допустимості:} параметризовані за класом дії, з більш високими порогами для змін индексу цінності, надзвичайних повноважень та інших критичних дій.
\item \textbf{Делегація:} Факультативна делегація з державними мандатами, розкриття конфліктів та відкликання.
\item \textbf{Перекращачі схеми:} надзвичайні паузи з зазначеними критеріями, уповноваженими операторами, умовами продовження роботи та необхідними постмортми.
\item \textbf{Прозорість:} опубліковані зміни та потоки, з окремим доказом розрахунку свапів, виконання емітентом, резервів, обмеження використання, маршрутизації та гарантів.
\end{itemize}
\textbf{Управління реєстрацією.} Розгортання CPP-сумістю може підтримувати реєстри відкриттів для ваучерів, токенів та Пулів. Авторизовані контролі можуть додавати, оновлювати, призупиняти або видаляти записи реєстру через розкритий процес управління розгортання. видалення реєстру впливає на виявлення і маршрутизацію через цей реєстр; не видаляє сам по собі токен, не змінює баланс власника, не виконує зобов'язання емітента або не відключає контракт, який іншим чином функціонує.

Опубліковані правила реєстру повинні передбачати умову статусу і можуть визначити повторну невиконання, шахрайство або неправильне представлення, небезпечну поведінку договору або постійне порушення опублікованих принципів як підстави для призупинення або видалення. Якщо це можливо, процес повинен передбачати повідомлення, можливість виправдання і шлях апеляції. Для аварійного виведення слід подати публічний звіт про інцидент і автоматично перевіряти або заходить сонце.

\textbf{Забороні списки.} Згідно з цим шаблоном реєстр не допускає:

\begin{enumerate}[leftmargin=*]
\item інструменти, які безпосередньо фінансують або стимулюють екологічне знищення за узгоджених кордонів, насильство або озброєння, примусову видобутку або системне зловживання; або
\item клас ваучерів, який не має чітких умов представлення і виконання, відповідальності та шляхів до виправдання.
\end{enumerate}
Список заборонених дій матиме версії, буде доступний для публічної перевірки й змінюватиметься лише через ухвалений поріг критичних дій і часову затримку, позначені як Q3 + T3 у Додатку D.

\subsection*{11.1 Обмінний курс та обмеження управління}
\addcontentsline{toc}{subsection}{11.1 Обмінний курс та обмеження управління}

\textbf{Зміни, зафіксовані в часі.} Розповсюдження за цим шаблоном змінило б методи обмінного курсу і окремо впроваджували параметри обмеження лише після публічного розкладу часу. На аварійній дорозі використовується окремо розкритий процес авторизації і включається автоматичний заход сонця або перегляд.

\textbf{Одобрений порог.} У шаблоні пропонуються більш високі порожники затвердження для змін бази индексу цінності та глобальних змін граничного рівня, середні порожні величини для змін сторонніх осіб, що стосуються конкретних фондів, та стандартні порожні величини для рутинних змін зборів.

\textbf{Опубліковані фільми.} Учасницький розподіл публікував би для кожної групи змінні індексу на ланцюжку, джерела або медиани оракулу, каденцію оновлення, обмеження вікнів і капелів, а також режими провалу або безпечні константи.

\textbf{Критерії надзвичайної перерви.} Учасницький розгортання передбачає умови, такі як розрив оракулу, високий ліміт використання в поєднанні з невдачами виконання, або невдачі інвариантів, разом з перевірками резюме і вимогами до перегляду після інциденту.

\subsection*{Приклад публічного індексового подачі для одного фонду та ваучера}
\addcontentsline{toc}{subsection}{Приклад публічного індексового подачі для одного фонду та ваучера}

\begin{itemize}[leftmargin=*]
\item \textbf{Символ:} Наприклад, \texttt{Maize\_50kg@IssuerY}.
\item \textbf{Референтна одиниця:} Індексний одиниць (IUX).
\item \textbf{Опублікована вартість:} 30.000 IUX.
\item \textbf{Джерело:} середня кількість ідентифікованих джерел, таких як дослідження місцевого ринку, бюлетень міністерства та базова лінія розгортання.
\item \textbf{Обновлення каденції:} щоденно в 18:00 EAT, з 24-годинним часовим блоком.
\item \textbf{Режим провалу:} заморожування на останньому дійсному значенні, застосування політики розкритого ліміту та зупинка після 72-годинної перерви.
\item \textbf{Розуміння:} опубліковані ноти та запис про зміни з попереднього оновлення.
\item \textbf{Підписанці:} розкриті багатозначні адреси та порог затвердження.
\end{itemize}
\subsection*{11.2 Пропозиційний план страхових фондів}
\addcontentsline{toc}{subsection}{11.2 Пропозиційний план страхових фондів}

\textbf{Вибиральний дизайн тільки.} Цей рулон застосовується лише до розпорядження, яке прямо прийняло та фінансувало страховий фонд та опублікувало покриті події, претенденти, відповідальний суб'єкт, активи, обмеження, виключення, вимоги до доказів, процес та умови управління. Ні поточний CLC App, ні GEF не забезпечують охоплення лише тому, що цей дизайн з'являється в Білої книзі.

\textbf{Можливі тригги.} Прийнята політика може охоплювати невиконання визначеної позиції емітента, дефіцит резервів групи або втрату моста або гарантії. технічний інцидент не має автоматичного призначення; відповідна політика контролює.

\textbf{Оцінка.} Відповідний орган згоджував би квитки з транзакціями, баланси запасів, облігації гаранта, записи представлення виплати, відповіді емітента та інші необхідні докази, а потім публікував запис інцидентів, який відповідає конфіденційності та законодавству.

\textbf{Ілюстративний водопад втрат.} Якщо кожен шар існує і законно застосовується, політика може використовувати: (1) відповідальні облігації емітента або ставки гаранта $\to$ (2) резерви на рівні фонду $\to$ (3) запропонований страховий фонд мережі $\to$ (4) тимчасове скорочення до вимоги до факультативного покриття, тільки якщо існуючі умови та чинне законодавство виразно дозволяють це $\to$ (5) законне відновлення за доведеного шахрайства або зловживання.

Корекція покриття не зменшує базовий зобов'язання видавниця щодо ваучера або не змінює баланс на ланцюзі, якщо не дозволяють виразно цей результат чинні вже існуючі умови та чинне законодавство, а також якщо не отримано будь-яка необхідна згода власника.

\textbf{Ограничення та виключення.} Опублікована охоплення визначатиме обмеження, допустимі представлення, докази, вікна претензій, виключені маршрути або події, географічні обмеження та ставлення до вичерпаних резервів. Виплата може бути нульовою після досягнення відповідних обмежень.

\textbf{Ілюстративний графік відновлення.} Якщо прийнято і опубліковано:

\begin{enumerate}[leftmargin=*]
\item вимоги були б отримані спочатку від відповідального емітента або гаранта облигації, потім від відповідних резервів резерву, а потім від запропонованого страхового фонду мережі;
\item будь-яке скорочення до вимоги щодо факультативного покриття обмежувалося б тим, що дозволяють попередні умови покриття та чинне законодавство, аж до опублікованої максимальної кількості інцидентів;
\item план відновлення може застосовувати зазначену частку відновленої вартості протягом зазначеного періоду, після чого будь-який залишився покритий дефіцит перетворюється на зареєстрований збитк у публічному послідовному обліку; і
\item кожна рішення надавала б квитанцію з інцидентом ID, пов'язаними вимогами і ваучерами, рішенням, планом відновлення і вікном апеляції.
\end{enumerate}
\subsection*{11.3 Рамка гарантів}
\addcontentsline{toc}{subsection}{11.3 Рамка гарантів}

У цьому розділі розрізняється відповідальність емітента, факультативні захисті фонду та гарантії третіх осіб. Пули можуть конкурувати за курацією, умовами та чітко запропонованою охороною, не передбачаючи, що CLC App, CPP, GEF або будь-яка ширша мережа автоматично гарантує ваучер.

\subsection*{Відповідальність базового видавниця}
\addcontentsline{toc}{subsection}{Відповідальність базового видавниця}

\begin{itemize}[leftmargin=*]
\item Кожен ваучер є, перш за все, відповідальністю його емітента. Емітент зобов'язується надати зазначений товар, послугу або законний грошовий еквівалент за його опублікованими умовами.
\item Видавці публікували б, хто може представити ваучер, що означає виконання, де і коли він доступний, які докази потрібні і які засоби захисту застосовуються.
\item Якщо емітент не виконує вимоги, основним відповідальним є емітент. Захищення Пулів або мереж застосовується тільки тоді, коли вони приймаються, фінансуються та розкриються окремо.
\end{itemize}
\subsection*{Факультативні захисті Пулів}
\addcontentsline{toc}{subsection}{Факультативні захисті Пулів}

Куратор Пулу може вирішити додати вузько визначену охорону до прийнятих ваучерів. Це не автоматично, і необхідно було б визначити відповідальну сторону, фінансування, допустимі події, ліпіки, вікна, докази, виключення та засоби захисту в метаданні групи та застосованих умовах.

Ілюстративні типи захисту включають:

\begin{enumerate}[leftmargin=*]
\item \textbf{Покриття резервних активів:} після підтвердженої невиказання видавником, відповідальний суб'єкт групи платить визначену суму в визначеному резервному активі з урахуванням його опублікованої галочки та доступних фінансованих резервів.
\item \textbf{Окно відміни:} після кваліфікаційного заходу, фонд пропонує обмежений часом маршрут обміну до попереднього або іншого затвердженого активу, за умови обмеження та запасу. Це захисність ліквідності, яка залежить від запасів, а не обіцянка того, що кожен обмін є зваротним.
\item \textbf{Альтернативне виконання:} відповідальна сторона організує затвердженого замінного постачальника в межах опублікованої величини або величини.
\item \textbf{Захист курсового діапазону:} для вибраних класів ваучерів група пропонує лише коригування покриття або компенсацію, зазначену в її існуючих умовах. Це не зменшує зобов'язання емітента щодо базового ваучера.
\end{enumerate}
\subsection*{Можливі джерела фінансування}
\addcontentsline{toc}{subsection}{Можливі джерела фінансування}

\begin{itemize}[leftmargin=*]
\item \textbf{Облигація емітента:} гарантії, що розміщені емітентом або зберігаються в розкритому резерві і доступні після перевірених покритих подій.
\item \textbf{Резерв Пулу:} активів, які контролюються відповідальним суб'єктом групи та виділяються на захисті, які він рекламує.
\item \textbf{Облигації гаранта третьої сторони:} обезпечення, розміщеного ідентифікованим зовнішнім гарантом для зазначених емітентів, класів ваучерів або подій.
\end{itemize}
У участі гаранта будуть дотримуватися опублікованих критеріїв допустимості, розміру облигацій, обмежень концентрації, повноважень до прийняття рішень та правил правопорядку.

\subsection*{Процес вимог}
\addcontentsline{toc}{subsection}{Процес вимог}

Прийнята політика визначала б триггерів, які підлягають аудиту, такі як термін виконання, який не виконується після дійсного представлення погашення, перевірена неплатіжність видавниця, покритий мост або нездахоплення гарантів або офіційно оголошений статус інцидентів. Він також визначав би:

\begin{itemize}[leftmargin=*]
\item як учасник відкриває претензію і надає необхідні докази представлення та виконання;
\item що перевіряє умови ваучера, відповіді емітента та технічні записи;
\item рішення та оскарження; і
\item Авторизований шлях виплати, активи, ліпені та квитанція.
\end{itemize}
Доходи від відновлення від емітентів, арбітражу або законного правопорядку будуть поповнювати відповідні облігації або резерви відповідно до опублікованої політики, перш ніж вони будуть використані для запропонованого доступу до обмінної мережі CLC.

\subsection*{Необхідні розкриття}
\addcontentsline{toc}{subsection}{Необхідні розкриття}

Для кожного покритого класу Пулів та ваучерів відповідальна сторона публікує:

\begin{itemize}[leftmargin=*]
\item чи відсутній гарант, не обов'язковий або обов'язковий;
\item розмір облигацій або резервів та обмеження концентрації;
\item типів захисту, активів, галовок, вікон і виключень;
\item терміни подачі, виконання, претензій та апеляції; і
\item ясна заява про те, хто гарантує що, а що не гарантовано.
\end{itemize}
\textbf{Принцип курації.} Куратори Пулів і відповідальні юридичні структури або структури управління несуть відповідальність за захист, який вони рекламують. Розгортання, що відповідає CPP-, може забезпечити стандарти, реєстри або факультативні спільні політики, але ні CLC, ні GEF автоматично не гарантують ваучери або Пули.

\subsection*{11.4 Охоронні огради проти захоплення}
\addcontentsline{toc}{subsection}{11.4 Охоронні огради проти захоплення}

В рамках цього шаблону наступні критичні дії вимагають прийнятого найвищого рівня затвердження та довгого періоду часу:

\begin{enumerate}[leftmargin=*]
\item зміна запропонованого платіжного водопаду, включаючи його охоплення та пріоритети основних операцій;
\item зміна корень канонічного реєстру;
\item зміна сфери покриття, обмеження вимог або влади до прийняття рішень;
\item розширення повноважень для аварійної перерви; або
\item Зневажаючи на те, що в цьому документі зазначені зобов'язання щодо форкабельності, прозорості або суверенітету групи.
\end{enumerate}
\subsection*{11.5 Процедура виходу і виходу}
\addcontentsline{toc}{subsection}{11.5 Процедура виходу і виходу}

Якщо управління було зафіксовано або цінності значно зменшилися, спільноти, Куратори Пулів та оператори могли б прагнути вийти за допомогою вилучення рівня управління мережами. Продовність базових пулів та ваучерів залежатиме від реалізованих контрактів, ключів, інтерфейсів, інфраструктури, послуг третіх осіб та відповідних зобов'язань.

Процес виходу може:

\begin{enumerate}[leftmargin=*]
\item \textbf{Опублікуйте стрічку:} експортують вибрані реєстри, ваучери, значення, обмеження та політику зборів, а потім публікують підписану хеш-шоу.
\item \textbf{Перерозподілення служб управління:} розгорнути нові коріння реєстру, послуги маршруту та будь-які прийняті модулі зборів або покриття в рамках нової відповідальної структури.
\item \textbf{Повторна реєстрація:} дозволяти Куратори Пулів прийняти участь, реєструючи свої адреси групи під новою корінням, не вимагаючи від власників міграції інакше функціональних ваучерів.
\item \textbf{Відновлення клієнтів:} додавати новий корінь як вибірковий профіль мережі в SDK та інтерфейсах, при будь-якій певній зміни, зроблені через розкритий процес управління.
\item \textbf{Управління мостовим періодом:} підтримувати сумісні маршрути, де вони безпечні, і відмовлятися від маршрутів, які порушують правила нового профілю.
\end{enumerate}
Метою проекту є те, що залишення канонічного реєстру не виключає функціональні локальні Пули. Фактична безперервність залишається залежною від розгортання; Федерація є оптимальним пластом відкриття та координації.


\section*{12. Пропозиційний план терміну програми ліквідності (незалежний облік)}
\addcontentsline{toc}{section}{12. Пропозиційний план терміну програми ліквідності (незалежний облік)}

Це ілюстративний перелік для майбутньої програми, окремо документованої. Це не пропозиція, актуальна функція додатку, або обіцянка ліквідності, доступу, страхування, повернення, вартості або виходу.

\begin{itemize}[leftmargin=*]
\item \textbf{Додаткові внески:} стабільні кошти, резервні токенів або затверджені ваучерів спільноти, як це зазначено в програмі.
\item \textbf{Місце розміщення:} назначається Пули зобов'язань або запропонований мережний фонд CLC на підставі окремо прийнятого мандата.
\item \textbf{Заключення і вихід:} спеціальні терміни програми, такі як перехідні періоди 30--90 днів та розкриті ворота під стресом.
\item \textbf{Кредит з платіжними платами:} факультативний механізм екс-пост під опублікованими оголомленнями і вікнами, а не обіцяна повернення.
\item \textbf{Розділ ризиків:} будь-яке скорочення до додаткового покриття повинно бути прямо дозволено існуючими умовами і чинним законодавством. Вона сама не змінює зобов'язання емітента з ваучера або баланс на ланцюзі.
\item \textbf{Повідомлення:} періодичні панелі та атестації, що охоплюють здоров'я Пулу, запас, обмеження, збори та будь-які фінансовані резерви покриття.
\item \textbf{Заповіді:} обмеження щодо використання доходів, контрольні дії оракулу, обмеження рівня, правила конфліктів та засоби захисту.
\item \textbf{Можлива допустимість запропонованого токану:} вносичі коштів, які мають призначені Пули, можуть мати право на отримання запропонованих грантів CLC на основі вимірюваного доступу до Пулу та підтвердженої виконання вимог емітента, за умови прийняття прийнятих умов, правил проти ігор та обмеження політики.
\end{itemize}


\section*{13. Глосар на простому мові}
\addcontentsline{toc}{section}{13. Глосар на простому мові}

\begin{itemize}[leftmargin=*]
\item \textbf{Cosmo-Local Credit (CLC):} Семейство продуктів, включаючи CLC App, документацію та більш широку роботу з об'єднання зобов'язань.
\item \textbf{CLC App:} Прогресивний веб-ап, який працює GEF на \texttt{cosmolocal.credit}.
\item \textbf{Протокол об'єднання зобов'язань (CPP):} Повторно використовувана модель курації, оцінки, обмеження, обміну та відповідального управління.
\item \textbf{Protocol v1.1.0:} Перевірений публічний випуск смарт-контракту.
\item \textbf{Пул зобов'язань:} Регулюється угода, яка визнає підтримувані активи і публікує правила обміну.
\item \textbf{\texttt{SwapPool}:} Сучасний контракт та контракт прямого обміну.
\item \textbf{Знак:} Баланс на ланцюзі або цифровий актив. Механіка токенів сама не створює викупної обіцянки.
\item \textbf{Ваучер:} Токен або запис, представлений емітентом як погашається зобов'язання за опублікованими умовами.
\item \textbf{Оферта:} Каталог товарів або послуг, пов'язаних з ваучером.
\item \textbf{емітент:} Сторона, відповідальна за публікацію та виконання зобов'язання про ваучер.
\item \textbf{Куратор Пулу:} Відповідальна структура, яка публікує і керує правилами Пул.
\item \textbf{Обмінний курс або котирування:} параметр транзакції, вироблений конфігурованим котированим знаком; не є доказом грошової суми або вартості погашення.
\item \textbf{Верхня межа балансу токена:} Поточний максимум \texttt{Limiter} для токенів, що зберігаються в Пулі. У Застосунку його можуть називати \textbf{«Кредитний ліміт»}.
\item \textbf{Обмін Пулу:} Директен обмін активами через один \texttt{SwapPool}.
\item \textbf{Виплата обміну:} Трансфери на ланцюгу та обліковий облік зборів, що завершують обмін групи.
\item \textbf{Представлення виплати:} Повернення або іншим чином представлення ваучерних одиниць емітенту.
\item \textbf{Виконання:} Емітент забезпечує обіцянений товар, послугу, вигоду або іншу продукцію.
\item \textbf{Випускання:} Запис, який запобігає повторному представленню виконаних одиниць.
\item \textbf{Пропанований маршрутизатор виконання:} Майбутнє програмне забезпечення, яке дозволить і виконає маршрути. Тільки поточні цитати \texttt{SwapRouter}.
\item \textbf{Пропанова CLC мережний Пул:} Майбутні домовленості щодо клірингу і координації бюджету на рівні мережі.
\item \textbf{запропонований токен управління CLC:} Будучий дизайн управління, а не поточний App або актив Protocol v1.1.0.
\item \textbf{Мережевий рейк:} Пропонована частка комісій групи передається на майбутній бюджет мережі.
\item \textbf{Гарантія або страхування:} Захист, який застосовується лише тоді, коли ідентифікована сторона прямо приймає, фінансує і публікує його.
\end{itemize}

\section*{14. Пропановані КПІ і показники здоров'я}
\addcontentsline{toc}{section}{14. Пропановані КПІ і показники здоров'я}

Розгортання повинно публікувати KPI тільки тоді, коли воно може визначити події, джерело, коорту або період, одиницю, метод оцінки та часовий штамп, виключення, політику виправлення, докази поза ланцюгом та обмеження якості даних.

До запропонованих заходів відносяться:

\begin{itemize}[leftmargin=*]
\item дійсні представлення виплати;
\item рівень виконання на основі кохорти;
\item повність списання;
\item латенція від представлення до виконання;
\item тривалість зберігання від придбання до подачі;
\item висутні допустимі зобов'язання за зазначеною методологією;
\item завершений обсяг обмінного фонду;
\item вимірений запас Пулу;
\item Використання балансового капіталу токенів фонду;
\item швидкість пропуску котирування, що зберігається окремо від швидкістю виконання маршруту;
\item допустимі резерви, поділені на експозиції, які чітко покриті;
\item вимоги гаранта та повернення в рамках ідентифікованої політики;
\item запропоновані плати за мережу та послуги, які були фактично отримані, без подвійного розрахунку брутних плати групи;
\item час виявлення управління, прийняття рішень, зупинки, ремонту та закриття;
\item скарги, суперечки, виправлення та засоби захисту; і
\item окремо підтверджуються соціальні або екологічні результати.
\end{itemize}
Дані про транзакції на ланцюжку самі по собі не підтверджують ідентичність, продуктивність емітента, задоволення, причинність, здоров'я спільноти або соціальний вплив. Ці твердження вимагають власної методології та доказів.

Ви бачите \href{https://docs.cosmolocal.credit/uk/white-paper/appendix-c-kpi-definitions}{Додатку В} для запропонованої специфікації вимірювання.


\section*{15. Дорожна карта (показательна)}
\addcontentsline{toc}{section}{15. Дорожна карта (показательна)}

\subsection*{15.1 Сучасний фундамент}
\addcontentsline{toc}{subsection}{15.1 Сучасний фундамент}

GEF управляє CLC App на \texttt{cosmolocal.credit} на ланцюжку Gnosis. Застосунок надає доступ до підтримуваних рахунків і гаманці, публічний каталог ринку та інтерфейси для токенів, ваучерів, пропозицій, Пули зобов'язань, переказів і прямих обмінів Пул.

Protocol v1.1.0 забезпечує основу поточного контракту: \texttt{GiftableToken}, безпосереднє виконання \texttt{SwapPool}, факультативні реєстри, модулі оцінки, обмеження балансу токенів Пул, плати Пул, додатковий протокольний honorarium \texttt{SwapRouter}. Доступність все ще залежить від інтерфейсу, розгорнутих контрактів, запасу, конфігурації, стану мережі, допустимості, постачальників, юрисдикції та опублікованих умов видавниця або групи.

\subsection*{15.2 Пропановані етапи}
\addcontentsline{toc}{subsection}{15.2 Пропановані етапи}

Ці названі етапи - це напрямки проектування, а не номера випуску, дати доставки або зобов'язання, які функція запустить.

\begin{itemize}[leftmargin=*]
\item \textbf{Фонд управління:} запропонований токен управління CLC; запропонований мережний Пул CLC; адаптери для оплати; Кворум, термін та основи управління.
\item \textbf{Маршрутизація і спостереження:} маршрутизатор SDK і реєстрні API; панелі охорони здоров'я; запропоновану основу страхової політики; наміри відновлення рівноваги; прототип партійної мережі.
\item \textbf{Інструменти ризику між доменами:} HTLC або маршрутизація на депозит; модулі гарантів, які регулюються окремо; заздалегідь встановлені регуляторні, розрахункові та рівнівні обмеження.
\item \textbf{Регулюваний доступ:} платіжні послуги третіх осіб, що відповідають конкретному призначенню; персональні мікробашки; виявлення послуг відповідності; аудити третіх осіб класів ваучерів.
\end{itemize}
Будь-яка платіжна послуга надається окремо визначеними третіми особами відповідно до чинної юрисдикції, допустимості, зборів, обмежень та умов. Контракти CLC App та Protocol v1.1.0 самі не працюють на фіатних рельсах.

Multi-hop виконання, спільне страхування, запропонований CLC управління токен і запропонований CLC мережний пуль, партійна мережа, персональні мікропули і регулювані платежні послуги залишаються запропонованими або залежно від впровадження, поки імплементація і його правила не визначить їх активними.


\section*{16. Пропозиційний шаблон оцінки}
\addcontentsline{toc}{section}{16. Пропозиційний шаблон оцінки}

Реестр, Пул або майбутня програма ліквідності можуть адаптувати цей шаблон. Це не є нинішнім універсальним стандартом або гарантією.

\begin{enumerate}[leftmargin=*]
\item \textbf{Відзначані факти:} Ідентифікація і історія емітента, умови ваучера, Оферти, докази потужності, конфігурація фонду, контролер, запас, обмеження, резерви, гаранти, інциденти і нерешенні зобов'язання.
\item \textbf{Мета та зацікавлені сторони:} намірена користь, ризики, конфлікти, екологічні та соціальні наслідки, а також хто несе втрати або відповідальність.
\item \textbf{Оцінка:} Офертна ціна, зазначена вартість ваучера, метод обмінного курсу групи, посилання на Oracle, одиниці, часові марки та причини будь-яких розбіжностей.
\item \textbf{Границі:} допустимість, заборонене використання, концентрація, обмеження на баланс токенів, паузи, географічні або юридичні обмеження та триггерів перегляду.
\item \textbf{Захисті:} ідентифікована зобов'язана особа, охоплена подія, фінансування, ліміт, виключення, тривалість, докази, процес подання претензій та засоби захисту.
\item \textbf{Рішення:} затверджують, переглядають, відкидають, призупиняють або збільшують їх для перегляду людьми; ідентифікують приймача рішень, причини, умови, дату перегляду та процес апеляції або виправлення.
\end{enumerate}
Вихід не повинен описувати курацію, обмеження, резерви або перевірку як підтвердження, потенціал емітента, страхування або гарантоване виконання, якщо відповідальна сторона явно не встановила цю охорону.


\section*{17. Правовий лист та запис про відповідність}
\addcontentsline{toc}{section}{17. Правовий лист та запис про відповідність}

CPP може координувати токенів, ваучерів, пулів та обмінних операцій, чий правовий режим залежить від їх розробки, маркетингу, використання, відповідальних сторін та юрисдикції. Ніщо в цій Білої книзі не є юридичною, податковою, інвестиційною, кредитною або фінансовою консультацією.

Використання CLC App регулюється \href{https://docs.cosmolocal.credit/uk/governance/terms}{Умови користування} GEF управляє додатком та підтримкою інфраструктури. Якщо GEF прямо не відіграє іншу роль, вона не є емітентом, Куратор Пулу, захисником, кредитором, позичальником, брокером, гарантом, викупником, радником, страховником або стороною зобов'язань між учасниками.

Залежні від розгортання платні послуги повинні визначити їх відповідального постачальника, юрисдикції, допустимість, збори, обмеження, модель зберігання та умови. Наявність підключачача не означає, що GEF або Пункт експлуатують основний регулюваний сервіс.

\subsection*{17.1 Розкриття ваучерів та пропозицій}
\addcontentsline{toc}{subsection}{17.1 Розкриття ваучерів та пропозицій}

Емітент повинен публікувати і зберігати актуальність:

\begin{itemize}[leftmargin=*]
\item відповідальна ідентифікація та контактна інформація;
\item відповідну пропозицію, одиницю, постачання, зазначену вартість та потужність;
\item місця представлення та процедури;
\item терміни виконання та докази;
\item термін дії, збори, податки, обмеження та правила змін матеріалу;
\item скарги, заміни, реституції або інші засоби захисту; і
\item метод викиду, який запобігає повторному використанню після виконання.
\end{itemize}
Токенний контракт не встановлює ці умови або не доказує виконання.

\subsection*{17.2 Розкриття Пулів}
\addcontentsline{toc}{subsection}{17.2 Розкриття Пулів}

У Куратор Пулу слід опублікувати:

\begin{itemize}[leftmargin=*]
\item ціль фонду та відповідальних осіб, що приймають рішення;
\item власник \texttt{SwapPool}, прокси-адміністратор, контролер залежності та одержувачі плати;
\item допущені активи та правила призупинення або видалення;
\item методи оцінки, джерела котирування, збори, загальні межі балансу токенів, запаси та повноваження для виведення власником;
\item права внеску і виходу;
\item будь-яка резерва, гарантія, гарант, страховий поліс та правило розподілу збитків; і
\item Управління, модернізація, надзвичайні процеси, скарги, міграція та припинення.
\end{itemize}
Ідентифікація не є підтвердженням, оцінкою, страхуванням або гарантією GEF.

\subsection*{17.3 Акції та зобов'язання}
\addcontentsline{toc}{subsection}{17.3 Акції та зобов'язання}

Зовсім звичайний \textbf{Обмін Пулу} біржі підтримуються активами за показанною котировкою та транзакційними лімітами. Направлення активів не робить його кредитом, погашенням, виплатою видавником або реалізацією в реальному світі.

\textbf{Пред'явлення для погашення} повертає або пред'являє Емітенту одиниці Ваучера. \textbf{Виконання} --- це обіцяна дія Емітента. \textbf{Списання} фіксує виконані одиниці, щоб їх не можна було використати повторно. У Застосунку дія \textbf{«Погасити»} зараз готує передавання власнику токена; саме це передавання не підтверджує Виконання.

У додатках \textbf{«Вивести ваучер з обігу»} Акція - це зворотне видалення каталогу з списку. Вона не спалює баланси, не скасує вимоги або не виконує зобов'язання.

Будучий кредит або інший кредитний продукт потребує окремо представлених додаткових і транзакційних умов перед використанням. Жоден звичайний відправка, внесок, депозит в фонд, обмін в фонд, подача, виконання або виплата не створюють кредит лише через його напрямок або тип активів.

\subsection*{17.4 Захист, докази та місцеве законодавство}
\addcontentsline{toc}{subsection}{17.4 Захист, докази та місцеве законодавство}

Лиміт, резерв, запис в реєстрі, перелік додатків або блокчейн-транзакція не є автоматично гарантією або страховою політикою. Будь-яка охорона повинна визначати відповідальну сторону, охоплену подію, фінансування, обмеження, виключення, тривалість, докази, процес подання претензій та чинні умови.

Докази з публічної ланцюги можуть показати адреси, активи, суми, часові марки та події договору. Вона сама по собі не підтверджує ідентичність, потенціал емітента, виконання, задоволення, обговорення управління, юридичне списання або соціальний вплив.

Функції можуть бути обмежені або недоступні через закон, санкції, допустимість, стан договору, запас, обмеження, безпеку, юрисдикцію або послуги третіх осіб. Емітенти, Куратори Пулів, постачальники та учасники залишаються відповідальними за визначення та дотримання чинного законодавства.


\section*{18. Висновок}
\addcontentsline{toc}{section}{18. Висновок}

Cosmo-Local Credit пов'язує поточне застосування і Protocol v1.1.0 фундамент з більш широким запропонованим проектом для самостійно керуваного Пули зобов'язань. Сучасні прямі обмінні котирування, компоненти котирування та обмеження та публічні записи транзакцій можуть підтримувати відповідальний обмін, але вони не доказують виконання емітентом, не створюють прав вкладчиків або не гарантують вартість, ліквідності, погашення, страхування, правовий статус або захист від втрат.

Пропанований маршрутизатор виконання, торгова система, активів управління, клірингу мережі, програм ліквідності та спільного захисту, описані в цій статті, потребують окремого впровадження, фінансування, управління та опублікованих термінів. Проект спрямований на розширення оперативної взаємодії при збереженні ефективності емітента, управління фондом та відповідальності в реальному світі з ідентифікованими сторонами.


\section*{А. Пропанована модель оцінки та плати}
\addcontentsline{toc}{section}{А. Пропанована модель оцінки та плати}

Цей додаток визначає запропоновану систему вимірювання. Protocol v1.1.0 не зафіксує виконання видавником, списання в реальному світі або будь-яке поле даних, яке потрібно нижче.

\subsection*{Означення події та запасу}
\addcontentsline{toc}{subsection}{Означення події та запасу}

Для ваучерного класу \emph{j}, коорти або періоду \emph{t}, і розкритого методу оцінки \emph{m}:

\begin{itemize}[leftmargin=*]
\item \texttt{O\_\{j,t,m\}}: вартість невиплачених кваліфікованих зобов'язань на межі вимірювання.
\item \texttt{X\_\{j,t,m\}}: вартість завершених обмінів фонду протягом періоду.
\item \texttt{P\_\{j,t\}}: одиниці, які дійсно представлені для погашення емітенту.
\item \texttt{F\_\{j,t\}}: представлені одиниці з окремо засвідченим виконанням емітентом.
\item \texttt{G\_\{j,t\}}: виконані одиниці з записом викиду, що запобігає повторному використанню.
\end{itemize}
\texttt{O} не може бути виведено тільки з пропозиції токенів. Політика оцінки повинна визначити відповідального емітента і виключати, відповідно до необхідності, інвентаризацію, що знаходиться в руках емітента, спалені одиниці, закінчені одиниці, звільнені одиниці, випробувальні баланси, недоступні баланси та токенів, умови яких не створюють непотрібного зобов'язання третіх осіб.

Кожна оцінювана мера повинна публікувати одиницю, джерело, метод оцінки, часовий штамп та обробку дивергентних курсів курсу. Передача на ланцюжку може підтримувати \texttt{X} або докази представлення; вона сама не встановлює \texttt{F} або \texttt{G}.

\subsection*{Мережа виконання на основі коортів}
\addcontentsline{toc}{subsection}{Мережа виконання на основі коортів}

Для кохорти дійсних представлень про погашення:

\texttt{FulfillmentRate = FulfilledPresentments / ValidPresentments}

\texttt{DischargeCompleteness = DischargedFulfillments / FulfilledPresentments}

Використовуйте одну і ту ж закриту або зрілу кохорту в кожному чисельніку і номінаторі. Доповідь відхилена, відкликана, закінчена, оскаржена, частково виконана, виправлена і ще відкрита подачі окремо.

\texttt{FulfillmentLatency = median(t\_fulfilled - t\_presented)}

\texttt{HoldingDuration = median(t\_presented - t\_acquired)}

Запізність виконання вимірює сервіс видавниця після представлення. Тривалість утримання є окремою мірою і не повинна бути позначеною як затримка виплати.

\subsection*{Різні швидкості}
\addcontentsline{toc}{subsection}{Різні швидкості}

Пропозиційна швидкість списання зобов'язань може бути розрахована лише у тих випадках, коли \texttt{O} і виконана вартість використовують один і той же метод оцінки:

\[
V_{\mathrm{commit},t} = \frac{\mathrm{FulfilledValue}_t}{\mathrm{AverageOutstandingEligibleCommitmentValue}_t}
\]

Мережа обмінної активності групи є окремою:

\texttt{V\_swap,t = PoolSwapValue\_t / AverageMeasuredPoolInventoryValue\_t}

Жодна з цих цінностей не підтверджує соціальний вплив, потенціал емітента, рентабельність або грошова конвертованість.

\subsection*{Пропанова доходів від мережі}
\addcontentsline{toc}{subsection}{Пропанова доходів від мережі}

Дозвольте:

\begin{itemize}[leftmargin=*]
\item \texttt{PF\_t} - брутні збори по складу, отримані протягом періоду;
\item \texttt{NR\_t} є запропонованим мережним рейком, який фактично отриманий як розкрита частка цих зборів групи;
\item \texttt{RF\_t} - окремі запропоновані або фактично отримані плати за маршрутизацію або обслуговування; і
\item \texttt{$\chi$\_t} є вимірюваною часткою отриманих доходів, яка є допустимою і конвертованою для зазначеного використання в грошовій валюті після витрат і політичних обмежень.
\end{itemize}
З точки зору запропонованого бюджету мережі:

\texttt{ProposedNetworkRevenue\_t = NR\_t + RF\_t}

\texttt{CashUsableNetworkRevenue\_t = $\chi$\_t $\times$ (NR\_t + RF\_t)}

Не додавайте \texttt{PF\_t} до \texttt{NR\_t}: рейк є перерахуванням від брутних зборів Пул і в іншому випадку буде розраховано двічі. Натомість поточний Protocol v1.1.0 підтримує додаткову плату протоколу; його квитки повинні бути зареєстровані окремо від цієї запропонованої моделі.


\section*{Б. Ілюстративний водопад майбутнього збору}
\addcontentsline{toc}{section}{Б. Ілюстративний водопад майбутнього збору}

Це запропонований проект для майбутнього розгортання. Вона застосовується тільки в тому випадку, якщо вона реалізується, фінансується та приймається за допомогою опублікованого управління та умов учасника. Він не описує поточну плату Protocol v1.1.0, право на внесник, резерву, гарантію або договору страхування.

Нехай \texttt{F\_in} є доходами, які насправді отримують запропонований бюджет мережі протягом епохи: запропонований рейк мережі, окремі збори за маршрутизацію або послуги, і будь-які інші чітко визначені доходи. Грошові збори за Пули, які залишаються у Пулах, не включені.

Дохідні активи повинні бути класифіковані як:

\begin{itemize}[leftmargin=*]
\item \texttt{E\_cash}: активи, які можуть бути використані у грошовій валюті відповідно до прийнятої політики; або
\item \texttt{E\_kind}: ваучер у натуральному вигляді або інші активи, які не розглядаються як грошові конвертації.
\end{itemize}
Будь-яка політика конвертації потребує реєстрації активів і місця проведення, відповідальних органів, джерел ціни, обмежень просліплення, звітності та відповідних правових контролів.

Пропозиційний водопад може застосовувати отримані доходи у цьому порядку:

\begin{enumerate}[leftmargin=*]
\item \textbf{Мета покритих резервів:} фінансувати прийняту резервну ціль або страхову ціль, засновану лише на визначених об'єктах покриття, допустимих активах, виключеннях та правилах вимог.
\item \textbf{Основні операції:} фінансувати розкритий, обмежений операційний бюджет.
\item \textbf{Мандати щодо ліквідності:} Фонд затверджений, окремо регулюється Пункт або маршрутування програм.
\item \textbf{Останній бюджет:} розподіляти будь-яку залишену суму між операціями, програмами ліквідності та додатковою буферу під опублікованими ліпіками.
\end{enumerate}
Ціль резерву не є гарантією. Будь-яке страхування або гарантія повинна ідентифікувати зобов'язану сторону, охоплену подію, фінансування, обмеження, виключення, тривалість, докази, процес пошкодження та розподіл втрат.

\textbf{Охоронний рельс:} Призначення водопадів призначені для прийнятого покриття, операцій та послуг щодо ліквідності. Вони не можуть бути оформлені або виконані як операції підтримки цін.

Згідно з запропонованою моделлю, будь-які запропоновані токенів CLC управління, придбані за допомогою окремо включенної програми зовнішньої ліквідності, були б виведені або розміщені в розкритому безголосовому лаві. Цей механізм не використовується.


\section*{C. запропоновані визначення KPI}
\addcontentsline{toc}{section}{C. запропоновані визначення KPI}

Ці КПІ утворюють запропоновану специфікацію вимірювання, а не заяву про те, що поточна програма або протокол записують кожну необхідну подію.

Кожен опублікований KPI повинен включати:

\begin{itemize}[leftmargin=*]
\item відповідальний видавець і джерело даних;
\item визначення події або стану;
\item кохорт або період вимірювання;
\item одиниць і метод оцінки;
\item терміновий знак оцінки;
\item правила включення та виключення;
\item обробка часткових, суперечливих, закінчених, недоступних та виправлених записів;
\item необхідний доказ або атестація поза ланцюгом;
\item історія перегляду та обмеження якості даних; і
\item чи є результат на ланцюжку, повідомлений, засвідчений, самостійно перевірені або оцінені.
\end{itemize}
\begin{longtable}{P{0.31\linewidth} P{0.31\linewidth} P{0.31\linewidth}}
\toprule
KPI & Пропанова визначення & Необхідні докази та виключення \\
\midrule
\endhead
\textbf{Підтверджені пред'явлення} & Кількість або одиниці, прийняті Емітентом до процесу погашення за визначений період & ідентифікатор пред'явлення, Емітент, дозвіл Власника, сума, час і статус; виключити дублікати та недійсні запити \\
\textbf{Рівень виконання} & Виповниті дійсні представлення, поділені на дійсні представлення для однієї і тієї ж зрілої кохорти & окремі докази ефективності емітента; звітувати про відкриті, відхилені, суперечливі, часткові та виправлені справи \\
\textbf{Повністю розряду} & Заповниті представлення з записом про списання, поділене на заповниті представлення & Спілення, скасування, відключення запису або інший доказ не повторного використання, пов'язаний з виконанням \\
\textbf{Запізність виконання} & Середній і 90-й перцентіль часу подачі до виконання & Не замінюйте час зберігання від видачі до подачі або від придбання до подачі \\
\textbf{Тривалість тримання} & Медиан і розподіл часу від придбання до подачі & Визначити події придбання та виключити невідомі часи придбання \\
\textbf{Вистатні допустимі зобов'язання} & Заможні зобов'язання третіх осіб, що залишаються після визначених виключень & ідентифікація і умови видавництва; виключити відповідний запас емітентів, термін дії, спалювання, списання, випробування та токенів без зобов'язань \\
\textbf{Об'єм обміну Пулом} & Вартість виконаних прямих обмінних операцій групи за одним розкритим методом оцінки & Події розрахунку на ланцюжку, токенні одиниці, джерело курсу, час; не класифікують як виконання емітентом \\
\textbf{Інвентаризація Пулів} & Вимірювані підтримувані активи, які містяться в базі в зазначений час & Контрактні залишки, резерви на збори, недоступні активи, метод оцінки та повноваження для виведення власником \\
\textbf{Достатність резервів} & Кваліфіковані, доступні резервні активи, розділені на експозиції, які чітко покриваються & Політика покриття, допустимість активів, зберігання/контроль, зобов'язання, виключення та оцінка; не загальний податок токенів за замовчуванням \\
\textbf{Ограничення використання} & Вимірений баланс токенів фонду, поділений на його поточний конфігурований капітал & Limiter адреса, токен, базар, часовий штамп, зміни та періоди без лімітера \\
\textbf{Цитатний пропускний показник} & Успішні відповіді на цитати розділені на дійсні спроби цитати & Результат тільки за цитатами; не виконання маршруту \\
\textbf{Рівень виконання маршруту} & Завершені виконання в декількох хопах розділені на дійсні спроби виконання & Використовується лише для реалізованої системи виконання; доповідь за правилами розбіжності та атомності \\
\textbf{Виплата гаранта} & Виплачувана виплата, отримана поділена на сплачені покриті позики & Визначений гарант, політика вимог, терміни, витрати, суперечки та списання \\
\textbf{Пропанова доходів від мережі} & Пропанований мережний рейк отриманий плюс отримані окремі тарифи маршрутування/сервісу & Виключити брутні податки на Пули, що зберігаються у Пулах, і уникнути подрахунку два рази \\
\textbf{Своєчасність управління} & Час виявлення, прийняття рішень, зупинки, ремонту та закриття інциденту & Визначені годинники, відповідальні органи, надзвичайні повноваження, апеляції та відсутні події \\
\bottomrule
\end{longtable}

Заявлення про соціальний вплив вимагають окремої методології. Активітет блокчейна сам по собі не встановлює ідентичність, продуктивність емітента, задоволення, здоров'я спільноти, причинність або вплив.


\section*{Д. Пропановані параметри запуску}
\addcontentsline{toc}{section}{Д. Пропановані параметри запуску}

Кожне значення нижче є ілюстративним. Це не поточне замовчування додатку, настройка Protocol v1.1.0, гарантія, пропозиція або зобов'язання щодо доставки. У майбутньому розгортання потрібно було б прийняти, застосувати, спостерігати і опублікувати свої фактичні параметри.

\begin{itemize}[leftmargin=*]
\item \textbf{Кворумні рівні для запропонованого токену управління CLC:}
\begin{itemize}[leftmargin=*]
\item Рутин Q1: не менше 4\% кворуму і більше 50\% затвердження.
\item Q2 чутливість: не менше 10\% кворуму і не менше 60\% затвердження.
\item Q3 критична: щонайменше 20\% кворуму і щонайменше 66.7\% затвердження.
\end{itemize}
\item \textbf{Пропановані терміни:}
\begin{itemize}[leftmargin=*]
\item T1: 48 годин для дій Q1.
\item T2: 7 днів для дій Q2.
\item T3: 30 днів для дій Q3.
\end{itemize}
\item \textbf{Каденція епохи:} 7 Дні, включаючи будь-які прийняті голосування, публікацію бюджету та запропоновані вікна sCLC.
\item \textbf{Надзвичайна пауза:} невідкладно через ідентифікований орган з надзвичайних ситуацій, термін дії якого закінчується через 72 години, якщо він не ратифікується відповідно до прийнятого процесу.
\item \textbf{Пропанований мережний рейк:} 20\% зборів учасника групи за випадом, що обмежується політикою.
\item \textbf{Ілюстративний діапазон плати за базен:} 0\%--20\%, опубліковані кожним учасним фондом.
\item \textbf{Предпанований обмеження маршрутного або обслуговуючого збору:} 20 bps на маршрут.
\item \textbf{Пропозиційний показник скорочення мережі:} \texttt{rake\_rate\_p = pool\_fee\_p $\times$ rake\_share\_p}.
\item \textbf{Виплата податків:} класифікують отримані активи як грошово допустимі \texttt{E\_cash} або в натуральному вигляді \texttt{E\_kind} за опублікованим методом.
\item \textbf{Контроль конверсії:} реєстратори активів та місця проведення, відповідальні органи, джерела цін, вікна часу, обмеження відкладу, обмеження та звітування.
\item \textbf{Бюджетні можливості:} публікує запропонований бюджет доступу до зборів тільки після того, як прийняті покриті резерви та операційні цілі будуть виконані; бюджет може бути нульовим.
\item \textbf{Різки ризику:} не піддають одному класі активів ризику іншого класу без чіткого, інформованого заходу відповідно до чинного законодавства.
\item \textbf{Пропановані обмеження розрахунку та рахунку:} відмовлятися від несанкціонованої міжкласової діяльності і застосовувати прийняті обмеження.
\item \textbf{Факультативне скорочення покриття:} тільки в тих випадках, коли це дозволяють вже існуючі умови покриття, необхідна згода та чинне законодавство; вона сама по собі не знижує зобов'язання емітента щодо ваучера або баланс на ланцюзі.
\item \textbf{Контроль зовнішньої ліквідності, якщо це можливо окремо:} публікують масштаби місця проведення, метод вимірювання, триггерів, обмеження витрат і обсягу, контроль виконання, аварійні зупинки та квитки. Не представляйте програму як символічну підтримку.
\end{itemize}
Сучасні плати Protocol v1.1.0 Пункт, додаткові протокольні плати та обмеження балансу токенів Пункт залишаються регулюються реалізованими контрактами і конфігурацією, а не ці запропоновані значення.


\section*{Е. Працюючий приклад --- запропонований мережний рейк}
\addcontentsline{toc}{section}{Е. Працюючий приклад --- запропонований мережний рейк}

Цей приклад ілюструє запропоновану модель грейка. Це не додатковий розрахунк протоколу-плати, який використовується Protocol v1.1.0.

Припустимо:

\begin{itemize}[leftmargin=*]
\item участвуючий фонд взимає 2.00\% плату за фонд;
\item запропонований мережний рейк отримує 20\% від цієї плати групи; і
\item 25\% отриманого ракету є допустимим і конвертованим для зазначеного використання в грошовій валюті після витрат.
\end{itemize}
Ефективний запропонований показник мережі на маршрутовану вартість становить:

\texttt{rake\_rate = 2.00\% $\times$ 20\% = 0.40\% = 40 bps}

Частина грошових коштів, яка може бути використана за передбачуваним фактором допустимості на 25\%, складає:

\texttt{cash\_usable\_rate = 40 bps $\times$ 25\% = 10 bps}

Якщо запропонований бюджет мережі має щомісячний вимог до грошової номінації в розмірі \$150,000 і ніяких інших доходів, ілюстративна маршрутована вартість, необхідна при грошовій користуючій ставці в розмірі 10 б/с, становить:

\texttt{required\_routed\_value = \$150,000 / 0.001 = \$150,000,000 per month}

Цей розрахунк не включає в себе брутні збори групи, що зберігаються групими. Додавши ці збори до мережного рейку, можна було б подвоїти джерело рейку.

Реальний аналіз бюджету також повинен бути опублікований:

\begin{itemize}[leftmargin=*]
\item фактичні квитанції за рейк та плату за послуги;
\item допустимості активів та витрат на конверсію;
\item Участь у базі та обсяг маршруту;
\item прийняті резервні та операційні цілі;
\item втрати, суперечки, корекції та недоступні активи; і
\item Сценарії чутливості, а не обіцяного зростання або повернення.
\end{itemize}
Будь-який запропонований бюджет для доступу до зборів або програм ліквідності залишився б нижче його прийнятих пріоритетів і міг би бути нульовим.


\section*{Ф. Перелік даних}
\addcontentsline{toc}{section}{Ф. Перелік даних}

\begin{enumerate}[leftmargin=*]
\item \textbf{Історичні докази Sarafu Network}
\begin{itemize}[leftmargin=*]
\item Архів датирований обсяг свапів, визначення активного користувача та активів, інциденти, презентації погашення, старіння ваучерів, підтвердженої виконання, списання, заходи тривалості утримування, періоди, виключення та методологія окремо від поточних метрик Cosmo-Local Credit.
\item Зберегти \href{https://dune.com/grassrootseconomics/sarafu-network}{історична панель Dune Analytics}.
\end{itemize}
\item \textbf{Пропозиція доказів виконання зборів, якщо вона реалізована}
\begin{itemize}[leftmargin=*]
\item Докажіть, що будь-яка платна гака і реєстрні шлюзи інтегровані в виконаний шлях транзакції, а не виконані тільки інтерфейсом.
\item Доставити випробування, що свідчать про те, що відповідна служба виконання відкидає посадку, отримання якої не підтверджує збору зборів відповідно до прийнятої політики.
\item Опублікуйте тестувальні вектори, випадки провалу, допустимі активи та відповідні Пули, пов'язані з ідентифікованою версією реєстру.
\item Звязати відповідні \href{https://software.grassecon.org}{опис програмного забезпечення}.
\end{itemize}
\item \textbf{Гарантор і пакет покриття}
\begin{itemize}[leftmargin=*]
\item Опублікуйте допустимість, відповідальні сторони, фінансування, обмеження, премії, якщо це необхідно, триггерів, доказів, виключень, повноважень до прийняття рішень та прикладів спорів і апеляцій.
\item Включити вибірковий ваучер і розкриття об'єкта, що відрізняють відповідальність видавниця від будь-якої чітко запропонованої захисту об'єкта або гарантії третіх осіб.
\item Зберегти \href{https://sarafu.network/reports}{історичні Sarafu Network звіти архів} і \href{https://youtu.be/5yGaP31bvs0?si=fZ-AMTEXsmVR8Ulv}{представлення історичного року рецензування} як докази лінії, а не доказ поточного покриття або прийняття CLC.
\end{itemize}
\item \textbf{Правовий пакет та пакет відповідності}
\begin{itemize}[leftmargin=*]
\item Документ стратегії юрисдикції, класів ваучерів, політики грошового еквіваленту, KYC або варіантів атестації, геооценування, розкриття інформації про споживачів, контроль неправильного продажу, а також відмінність між звичайним обміном Пул та продуктом експресного кредитування.
\item Зв'язок ефективної \href{https://docs.cosmolocal.credit/uk/governance/terms}{Cosmo-Local Credit Умови надання послуг} і збереження замінених версій через контрольовані записи.
\end{itemize}
\item \textbf{Засилення управління}
\begin{itemize}[leftmargin=*]
\item Указати процедури конфлікту інтересів та відмовлення, реєстрацію рішень та апеляцій, санкції, правовий процес видалення реєстру, обмеження надзвичайних потужностей та приклади змін валютного курсу та обмежень, що зафіксовані вчасно.
\end{itemize}
\item \textbf{Витоки та упаковка гарантій}
\begin{itemize}[leftmargin=*]
\item Публікуйте перелік джерел і ліцензій, розгорнуті адреси і версії, доступні аудиторські звіти, інварианти, походження конструкції, повноваження адміністратора, контакти з інцидентами та панелі управління.
\item Звязати відповідні \href{https://github.com/grassrootseconomics}{Репозиторий "Грассройт Економіка"}.
\end{itemize}
\end{enumerate}

\end{document}
